\documentclass[a4paper,12pt]{article}
% \usepackage{arxiv}
\usepackage{xeCJK}
\usepackage{amsmath,amssymb,amsthm,stackrel,nccmath,bm}
\usepackage{color,xcolor}
\usepackage{fontspec}
\usepackage[lite,subscriptcorrection,slantedGreek]{mtpro2}
\linespread{1.5}
\usepackage{hyperref}
\hypersetup{
    final,                   % Include links even if document is in 'draft' mode.
    hidelinks,               % Prevent boxes from being drawn around links in some viewers.
    breaklinks=true,         % Allow line breaks in the middle of links
    colorlinks=true,
    linkcolor=black,         % TOC, links to labeled equations and environments
    urlcolor=black!30!blue,  % URLS including links in references
    anchorcolor=blue,        % I'm not sure what this does.
    citecolor=black!30!blue, % In-line citations  
}
\usepackage[
    noabbrev,               % E.g., "Figure" instead of "Fig."
    capitalise,             % E.g., "Figure" instead of "figure"
    nameinlink              % Make "Figure 1" linked instead of just "1".
]{cleveref}
\usepackage{cite}
% \usepackage{thmbox}
\setmainfont{Times New Roman}
\setCJKmainfont{SimSun}
\newtheoremstyle{sltheorem}
                {}          % Space above
                {}          % Space below
                {\upshape}  % Body font
                {}          % Indent amount
                {\bfseries} % Head font
                {:}         % Punctuation after head
                { }         % Space after theorem head
                {}          % Theorem head spec
% Enable the new "sltheorem" theorem style.
\theoremstyle{sltheorem}
\allowdisplaybreaks[4]
\makeatletter
\def\@endtheorem{\endtrivlist}
\makeatother
\newtheorem{definition}{Definition}[section]
\newtheorem{conjecture}{Conjecture}[section]
\newtheorem{proposition}{Proposition}[definition]
\newtheorem{theorem}{Theorem}[section]
\newtheorem{lemma}{Lemma}[theorem]
\newtheorem{corollary}{Corollary}[theorem]
\newtheorem{remark}{Remark}[theorem]
\newtheorem{example}{Example}[section]
\everymath{\displaystyle}

\DeclareMathOperator{\fmod}{mod}
\DeclareMathOperator{\val}{val}
\DeclareMathOperator{\sgn}{sgn}
\DeclareMathOperator{\ord}{ord}
\DeclareMathOperator{\fac}{fac}
\DeclareMathOperator{\rad}{rad}
\DeclareMathOperator{\raf}{raf}
\DeclareMathOperator{\ram}{ram}
\DeclareMathOperator{\Log}{Log}
\DeclareMathOperator{\rdu}{red}
\DeclareMathOperator*{\lcm}{lcm}

% \title{{On Dynamical Systems With $a_{k+1}\cdot\beta^{\Delta_{k+1}} = \alpha\cdot a_{k} + \gamma$ as Recursive Rule}}
\title{On the General Collatz Sequence: $\alpha \cdot a + \gamma$ reduced by $\beta$}
\author{Yaokai Liu
\href{mailto:yaokai.liu@foxmail.com}{{\texttt{yaokai.liu@foxmail.com}}}%
}
\begin{document}
\maketitle
\begin{abstract}
    This paper explores various properties of generalized Collatz sequences of the form $a_{k+1}\cdot\beta^{\val(a;k+1)} = \alpha\cdot a_{k} + \gamma$, extends this formulation to the field of rational numbers $\mathbb{Q}$, and constructs a hierarchical model for the set of source values. We define a specific density capable of preserving the latent algebraic structures among these sets, thereby mapping the discrete iterative process onto a continuous algebraic landscape.
    This paper proposes several propositions deemed crucial for proving the "$3n+1$ Conjecture" and offers an interpretation of the difference of powers problems from the perspective of dynamical systems.
    % And a proof of the $3n + 1$ conjecture is then provided.
\end{abstract}
\newpage
\tableofcontents
\newpage

% !TEX root = main.tex
\section{Introduction And the General Collatz Sequence}\label{sec:introduction-and-the-general-collatz-sequence} % (fold)
Our story starts from the $3n + 1$ Conjecture, hereinafter referred to as the Collatz conjecture, which is defined as follows:
\begin{conjecture}[Collatz Conjecture]~\label{conjecture:original-conjecture}
Given a number $a_0 = a \in \mathbb{N}$, we have the modular arithmetic notation
\[a_{k+1} = \left\{\begin{aligned}
    &a_{k}/2, &a_{k}\equiv 0\pmod{2}, \\
    &3a_{k}+1, &a_{k}\equiv 1\pmod{2},
\end{aligned}\right.\]
where $k \in \mathbb{N}$.
The conjecture is, for all $a_0 \in \mathbb{N}$, there's $k \ge 0$, that $a_{k} = 1$.
\end{conjecture}

First, for all $k$, there is $m \ge 0$ such that $3a_{k}+1 \equiv 0 \pmod{2^{m}}$.
We may wander the maximum value of $m$ which depends on $a_k$, what the $a_k$ depends on $a$ and $k$.
Let's define the symbol $\val_{2}(a;k+1)$ as count of factor $2$ in $3a_{k}+1$, where $a$ is the initial number given, and $k$ runs in $\mathbb{N}$.
With this we can redefine the sequence with
\[ a_{k + 1} = \frac{3a_{k}+1}{2^{\val_{2}(a;k+1)}}, \]
which ensures the $a_k$ always is an odd number.

It makes things easier.
Because when I start iteratively to apply the definition of $a_{k-1}$ into the definition of $a_{k}$, I found
\[ a_k \cdot 2^{\sum_{j=0}^{k}\val_{2}(a;j)} = a \cdot 3 ^ k + \sum_{i = 0}^{k - 1}2^{\sum_{j=0}^{i}\val_{2}(a;j)}\cdot3^{k-i-1}.  \]
Though it looks like a little bit complex, but with the definition $\mathcal{A}_{k} = \sum_{j=0}^{k}\val_{2}(a;j)$, the looks of the formula soon be better:
\[ a_k \cdot 2^{\mathcal{A}_k} = a \cdot 3 ^ k + \sum_{i = 0}^{k - 1}2^{\mathcal{A}_i}\cdot3^{k-i-1}. \label{equation:collatz-inductioon-3-2-1}\]
That means
\[ a \equiv -3^{-1} \cdot \sum_{i = 0}^{k - 1}2^{\mathcal{A}_i}\cdot3^{-i} \pmod{2^{\mathcal{A}_k}}. \]
Notice that $a$ is a positive finite number but $2^{\mathcal{A}_k}$ increase to infinity, that means $\sum_{i = 0}^{k - 1}2^{\mathcal{A}_i}\cdot3^{-i}$ should be a finite odd times of $3^{-1}$ under the modulo.
This congruence equation system has one obvious solution: $a = 1$, that $\mathcal{A}_k = 2k$, holds
\[ -3^{-1} \cdot \sum_{i = 0}^{k - 1}2^{\mathcal{A}_i}\cdot3^{-i} = - 3^{-k}\cdot 4^k + 1 \equiv 1 \pmod{2^{2k} = 4^{k}}, \]
follows the Collatz conjecture.
We can see that each number generates a sequence $\mathcal{A}_{k}$ such that the above congruence equation holds.
So it is the time to explore what the arguments of these equations has a stability solution.
\begin{definition}[Collatz Sequence and Collatz Accumulation]~\label{definition:collatz-sequence}
    Given three arguments $\alpha, \beta, \gamma \in \mathbb{N}_{\ge1}$ that $\gcd(\alpha, \beta) = 1$, and $\alpha \ne 1, \beta\ne 1$.\newline
    For number $a \in \DiffSet{\mathbb{N}}{\beta\mathbb{N}}$, define the sequence $\left\{a_{_k}\right\}_{k = 0}^{\infty} \subset \DiffSet{\mathbb{N}}{\beta\mathbb{N}}$ named Collatz sequence and $\val_{\beta}(a;k)\in\mathbb{N}$ named Collatz reduce with
    \begin{align*}
        a_0 &=a, \\
        a_{k+1} \cdot \beta^{\val_{\beta}(a;k + 1)} &= \alpha\cdot a_{k} + \gamma.
    \end{align*}
    For a better looking, denote 
    \[\mathcal{A}_n^{m} = \sum_{k=m}^{n-1}\val_{\beta}(a;k + 1)\]
    named Collatz accumulation from $m$ to $n$.
    Specially when $m = 0$, we omit the superscript, that $\mathcal{A}_n = \mathcal{A}_n^{0}$.
\end{definition}

To avoid repeating the explanation each time, unless otherwise specified, we will use the uppercase calligraphic font to represent the Collatz accumulation of the corresponding letters by default.
Actually, $\val_{\beta}(a;k+1)$ is the $\beta$-adic valuation of $\alpha\cdot a_{k} + \gamma$.
And we can see, it is well-defined, because there is
\begin{proposition}[Uniqueness]~\label{proposition:uniqueness}
If $a = b$, then $\forall k \in \mathbb{N}$, there is $a_k = b_k$ and $\mathcal{A}_k = \mathcal{B}_k$.
\end{proposition}
\begin{proof}
    First, we have
    \begin{align*}
        a = b &\Rightarrow a_0 = b_0, \\
        a_1 \cdot \beta^{\val_{\beta}(a;1)} &= b_1 \cdot \beta^{\val_{\beta}(b;1)}, \\
        \mathcal{A}_1 = \val_{\beta}(a;1) &= \val_{\beta}(b;1) = \mathcal{B}_1.
    \end{align*}
    Second, if $a_k = b_k$ and $\mathcal{A}_k = \mathcal{B}_k$, we have
    \[a_{k+1} \cdot \beta^{\val_{\beta}(a;k + 1)} = \alpha\cdot a_{_k} + \gamma = \alpha\cdot b_{_k} + \gamma = b_{k+1} \cdot \beta^{\val_{\beta}(b;k + 1)}.\]
    That means
    \[a_{k+1} = b_{k+1} \cdot \beta^{\val_{\beta}(b;k + 1) - \val_{\beta}(a;k + 1)}. \]
    But there is $a_{k+1} \in \mathbb{N} - \beta\mathbb{N}$ as definition that $\beta \nmid a_{k+1}$, which means
    \[\val_{\beta}(a;k + 1) = \val_{\beta}(b;k + 1), \]
    and then
    \begin{align*}
        a_{k+1} &= b_{k+1}, \\
        \mathcal{A}_{k+1} = \mathcal{A}_{k} + \val_{\beta}(a;k+1) &= \mathcal{B}_{k} + \val_{\beta}(b;k+1) = \mathcal{B}_{k+1}.
    \end{align*}
    Sum up, we proved this proposition.
\end{proof}
\begin{proposition}
    If $\exists t > 0$ that $a_{t} = a_0 = a$, 
    then $\forall k > 0$, and $s = 1, 2, 3, \cdots, t$, there are
    $a_{k\cdot t+s} = a_s$ and $\mathcal{A}_{k\cdot t + s} = \mathcal{A}_{s}$.
\end{proposition}
\begin{proposition}[Co-uniqueness]~\label{proposition:co-uniqueness}
    If there are two numbers $a, b \in \DiffSet{\mathbb{N}}{\beta\mathbb{N}}$ that $\exists n \ge 0$, such that $\forall m = 0, 1, 2, \cdots, n$, there is $\mathcal{A}_m = \mathcal{B}_m$, then $\beta^{\mathcal{A}_n} \mid (a - b)$ and $\alpha^{n} \mid (a_n - b_n)$.
\end{proposition}
% \begin{proof} Trivial with \cref{theorem:inductive-formula}. \end{proof}
\begin{proof}
    We will use \cref{theorem:inductive-formula} to prove this proposition.\newline
    Since $\mathcal{A}_m = \mathcal{B}_m$, we have $\mathfrak{A}_{m} = \mathfrak{B}_{m}$.
    With the \cref{theorem:inductive-formula}, we have
    \[ (a_{m} - b_{m}) \cdot \beta^{\mathcal{A}_m} = (a - b) \cdot \alpha^{m}. \]
    Since $\gcd(\alpha, \beta) = 1$, then $\gcd(\beta^{\mathcal{A}_m}, \alpha^{m}) = 1$.
    So we have
    \[\beta^{\mathcal{A}_m} \mid (a - b), \quad \alpha^{m} \mid (a_m - b_m) \]
    holds for all $m\le n$.
\end{proof}
\begin{proposition}
    If $\forall m \ge 0$, there is $\mathcal{A}_m = \mathcal{B}_m$, then $a = b$.
\end{proposition}
\begin{proposition}
    $\forall n \ge k \ge m$ there is $\mathcal{A}_n^m = \mathcal{A}_n^{k} + \mathcal{A}_k^{m}$.
\end{proposition}
\begin{proposition}~\label{proposition:trivial-periodic-sequence}
    If there is $\Delta \in \mathbb{N}_{\ge1}$ holds $(\beta^{\Delta} - \alpha) \mid \gamma$, then let $a_0 = \frac{\gamma}{\beta^{\Delta} - \alpha}$, we have $\forall k \in \mathbb{N}$, $a_k = a_0$ and $A_k = k\cdot \Delta$.
\end{proposition}
\begin{proposition}~\label{proposition:equivalent-between-collatz-conjecture-and-our-definition}
Let arguments $(\alpha, \beta, \gamma) = (3, 2, 1)$.
For all odd number $a$, let $a'_{i_k}$ be the $k$-th odd number in \cref{conjecture:original-conjecture} and $a_k$ be the $k$-th element of Collatz sequence of $a$, then $\forall k \in \mathbb{N}$,  we have $a_{k} = a'_{i_{k}}$.
\end{proposition}
% \begin{proof} Trivial. \end{proof}
\begin{proof}
    Let $a'_{i}$ be the $i$-th number in \cref{conjecture:original-conjecture}, and $a'_{i_k}$ be the $k$-th odd number in it.\newline
    First, we have $a_{0} = a = a'_0 = a'_{i_0}$;\newline
    Second, if $a_{k} = a'_{i_k}$, suppose $2^{s_{k+1}} \mid 3a'_{i_k} + 1$ and $2^{s_{k+1}+1} \nmid 3a'_{i_k} + 1$.
    That means
    \begin{align*}
        a'_{i_k+1} &= (3a'_{i_k} + 1) / 2,   \\
        a'_{i_k+2} &= (3a'_{i_k} + 1) / 2^2, \\
        a'_{i_k+3} &= (3a'_{i_k} + 1) / 2^3, \\
        &\cdots \\
        a'_{i_k+s_k} &= (3a'_{i_k} + 1) / 2^{s_{k+1}} = a'_{i_{k+1}}
    \end{align*}
    and $a_{k+1} = (3a_k + 1) / 2^{\val_{\beta}(a;k+1)} = (3a'_{i_k} + 1) / 2^{\val_{\beta}(a;k+1)}$.
    That means
    \[a_{k+1} = a'_{i_{k+1}} \cdot 2^{s_{k+1} - \val_{\beta}(a;k+1)}.\]
    Notice that $2\nmid a_{k+1}$, means
    \begin{align*}
        s_{k+1} &= \val_{\beta}(a;k+1), \\
        a_{k+1} &= a'_{i_{k+1}}.
    \end{align*}
    Sum up, we proved this proposition.
\end{proof}
\cref{proposition:equivalent-between-collatz-conjecture-and-our-definition} ensures our definition of Collatz sequence is equivalent with the iterative function in \cref{conjecture:original-conjecture}.
Because for even numbers, we can easily figure out that every even number will reduce to an odd number at first.
In that situation, the proposition holds as well.
With this proposition, we can describe the Collatz Conjecture with the Collatz sequence:

\begin{conjecture}[Rewritten Collatz Conjecture]~\label{definition:rewritten-conjecture}
With arguments $(\alpha, \beta, \gamma) = (3, 2, 1)$,
for any odd number $a \in \mathbb{N}$, it's Collatz sequence will finally reach to $1$.
\end{conjecture}

The following discussion involves many terms like $\textstyle \sum_{i = 0}^{k - 1}\beta^{\mathcal{A}_i}\cdot \alpha^{k-i-1}$.
For a better looking, we first define a notation to represent them:
\begin{definition}~\label{definition:collatz-complement}
    Given a sequence of number $\{v_{i}\}\subset\mathbb{N}$, finite or infinite terms. $\forall m, n \in \mathbb{N}$ that $n \ge m$, denote $\mathcal{A}_{n}^{m} = \sum_{i=m}^{n-1}v_{i}$, and define $\mathfrak{A}_n^{m}$ called Collatz complement of $\{v_{i}\}$ at $n$ relative to $m$ with
    \[ \mathfrak{A}_n^{m} = \sum_{i = m}^{n - 1}\beta^{\mathcal{A}_i^{m}}\cdot\alpha^{n-i-1} \in \mathbb{N}. \]
    And define $\mathfrak{a}_n^{m}$ called reduced Collatz complement of $a$ at $n$ relative to $m$ with
    \[ \mathfrak{a}_n^{m} = \frac{\mathfrak{A}_n^{m}}{\alpha^{n-1}} = \sum_{i = m}^{n - 1}\frac{\beta^{\mathcal{A}_i^{m}}}{\alpha^{i}} \in \mathbb{Q}. \]
    Also, they can be omitted the superscript $m$ when $m = 0$ that $\mathfrak{A}_n = \mathfrak{A}_n^{0}$, and $\mathfrak{a}_n = \mathfrak{a}_n^{0}$.
\end{definition}

To avoid repeating the explanation each time, without further explanation,
we will use the same letter of the Collatz accumulation but uppercase fraktur font to represent Collatz complement,
and the same letter but lowercase fraktur font to represent reduced Collatz complement.
And if $\{v_{i}\}$ is the Collatz reduces of a number, we also say that $\mathfrak{A}_{n}^{m}$ and $\mathfrak{a}_{n}^{m}$ are of that number.

For the Collatz complement, we have two basic propositions:
\begin{proposition}~\label{proposition:collatz-complement-basic-prop1}
$\mathfrak{A}_{n}^{n} = 0$, and $\mathfrak{A}_{n+1}^{n} = 1$.
And $\mathfrak{a}_{n}^{n} = 0$, and $\mathfrak{a}_{n+1}^{n} = \frac{1}{\alpha^{n}}$.
\end{proposition}
% \begin{proposition}
%     $\mathfrak{A}_{n+1} = \alpha \cdot \mathfrak{A}_{n} + \beta^{\mathcal{A}_n}$. And $\mathfrak{a}_{n+1} = \mathfrak{a}_{n} + \frac{\beta^{\mathcal{A}_n}}{\alpha^{n}}$.
% \end{proposition}
\begin{proposition}~\label{proposition:collatz-complement-basic-prop2}
    $\forall n \ge k \ge m$, we have
    \begin{align*}
        \mathfrak{A}_n^{m} &= \beta^{\mathcal{A}_k^{m}}\cdot\mathfrak{A}_n^{k} + \mathfrak{A}_k^m\cdot\alpha^{n-k}, \\
        \mathfrak{a}_n^{m} &= \beta^{\mathcal{A}_k^{m}}\cdot\mathfrak{a}_n^{k} + \mathfrak{a}_k^m. 
    \end{align*}
\end{proposition}
\begin{proof}
\begin{align*}
    \mathfrak{A}_n^{m} =&\sum_{i = m}^{n - 1}\beta^{\mathcal{A}_i^{m}}\cdot\alpha^{n-i-1} \\
    =&\ \sum_{i = k}^{n - 1}\beta^{\mathcal{A}_i^{m}}\cdot\alpha^{n-i-1} + \sum_{i = m}^{k - 1}\beta^{\mathcal{A}_i^{m}}\cdot\alpha^{n-i-1} \\
    =&\ \sum_{i = k}^{n - 1}\beta^{\mathcal{A}_k^{m}}\cdot\beta^{\mathcal{A}_i^{k}}\cdot\alpha^{n-i-1} + \sum_{i = m}^{k - 1}\beta^{\mathcal{A}_i^{m}}\cdot\alpha^{k-i-1}\cdot\alpha^{n-k} \\
    =&\ \beta^{\mathcal{A}_k^{m}}\cdot\sum_{i = k}^{n - 1}\beta^{\mathcal{A}_i^{k}}\cdot\alpha^{n-i-1} + \sum_{i = m}^{k - 1}\beta^{\mathcal{A}_i^{m}}\cdot\alpha^{k-i-1}\cdot\alpha^{n-k} \\
    =&\ \beta^{\mathcal{A}_k^{m}}\cdot\mathfrak{A}_n^{k} + \mathfrak{A}_k^m\cdot\alpha^{n-k}, \\
    \mathfrak{a}_n^{m} =&\sum_{i = m}^{n - 1}\frac{\beta^{\mathcal{A}_i^{m}}}{\alpha^{i}} \\
    =&\ \sum_{i = k}^{n - 1}\frac{\beta^{\mathcal{A}_i^{m}}}{\alpha^{i}} + \sum_{i = m}^{k - 1}\frac{\beta^{\mathcal{A}_i^{m}}}{\alpha^{i}} \\
    =&\ \beta^{\mathcal{A}_k^{m}}\cdot\sum_{i = k}^{n - 1}\frac{\beta^{\mathcal{A}_i^{k}}}{\alpha^{i}} + \sum_{i = m}^{k - 1}\frac{\beta^{\mathcal{A}_i^{m}}}{\alpha^{i}} \\
    =&\ \beta^{\mathcal{A}_k^{m}}\cdot\mathfrak{a}_n^k + \mathfrak{a}_k^m.
\end{align*}
\end{proof}
Actually, we can formally define the Collatz complement only from \cref{proposition:collatz-complement-basic-prop1} and \cref{proposition:collatz-complement-basic-prop2}. More essentially, the two propositions are just the basic rules of the quasi $(\sigma, \tau)$-derivation of \cref{sec:a-kind-of-derivation-on-groupoid}, which in other words, the Collatz complement is a kind of $(\sigma, \tau)$-derivation on groupoid (\cref{def:quasi-derivation-on-groupoid}).

According to \cref{proposition:uniqueness} and \cref{proposition:co-uniqueness}, we have
\begin{proposition}[Uniqueness]~\label{proposition:uniquess-of-Collatz-complement}
    If $a = b$, then $\forall m, n \ge 0$, there are $\mathfrak{A}_n^{m} = \mathfrak{B}_n^{m}$ and $\mathfrak{a}_n^{m} = \mathfrak{b}_n^{m}$.
\end{proposition}
\begin{proposition}
    If $\forall v_i \ge 1$, then $\beta\nmid \mathfrak{A}_n^{m}$ and $\alpha \nmid \mathfrak{A}_n^{m}$.
\end{proposition}

Now we are ready to explore the properties of Collatz sequences.
But before that, I need to point out that there are many symbols in number theory.
These symbols, while simple, differ slightly from those in some other papers.
These symbols are defined in \cref{sec:number-theory-tools}. Before proceeding, please ensure you are familiar with these symbols first.


% !TEX root = main.tex
\section{The basic formula of Collatz Sequence}\label{sec:the-basic-formula-of-collatz-sequence}

With the definition of Collatz complement, we can easily find that the equation in \cref{equation:collatz-inductioon-3-2-1} also hold with all general Collatz sequences:
\begin{theorem}~\label{theorem:inductive-formula}
$\forall a \in \widehat{\mathbb{N}}_{{\beta}}$, $\forall k \in \mathbb{N}$ we have
\[ a_k \cdot \beta^{\mathcal{A}_k} = a \cdot\alpha^k + \gamma\cdot\mathfrak{A}_{k}. \]
\end{theorem}
% \begin{proof} Trivial. \end{proof}
\begin{proof}
    First, we have $a_0 \cdot \beta^{\mathcal{A}_0} = a = a \cdot \alpha^{0} + \gamma \cdot \mathfrak{A}_0$ and $a_1 \cdot \beta^{\mathcal{A}_1} = \alpha \cdot a + \gamma = a \cdot \alpha^{1} + \gamma \cdot \mathfrak{A}_1$.\newline
    Second, if $a_k \cdot \beta^{\mathcal{A}_k} = a \cdot\alpha^k + \gamma\cdot\mathfrak{A}_{k}$, then
    \begin{align*}
        a_{k+1} \cdot \beta^{\mathcal{A}_{k+1}}
        =&\ a_{k+1} \cdot \beta^{\val(a;k+1)} \cdot \beta^{\mathcal{A}_{k}} \\
        =&\ \alpha \cdot a_k \cdot \beta^{\mathcal{A}_k} + \gamma \cdot \beta^{\mathcal{A}_k} \\
        % =&\ \alpha \cdot (a \cdot\alpha^k + \gamma\cdot\mathfrak{A}_{k}) + \gamma \cdot \beta^{\mathcal{A}_k} \\
        =&\ a \cdot \alpha^{k+1} + \gamma \cdot \alpha \cdot \mathfrak{A}_{k} + \gamma \cdot \beta^{\mathcal{A}_k} \\
        % =&\ a \cdot \alpha^{k+1} + \gamma \cdot (\alpha \cdot \mathfrak{A}_{k} + \cdot \beta^{\mathcal{A}_k}) \\
        =&\ a \cdot \alpha^{k+1} + \gamma \cdot \mathfrak{A}_{k+1}.
    \end{align*}
    Sum up, we proved this theorem.
\end{proof}
It is very important and fundamental, and we will see later that all types of Collatz sequences are actually some kind of variant expression of the \cref{theorem:inductive-formula}. And this formula also explains why we call $\mathfrak{A}_k$ as the Collatz complement.
\begin{corollary}
$\forall m, n \ge 0$ we have
\[ a_{n+m} \cdot \beta^{\mathcal{A}_{n+m}^{m}} = a_{m} \cdot\alpha^{n} + \gamma\cdot\mathfrak{A}_{n+m}^{m}. \]
\end{corollary}
\begin{corollary}~\label{corollary:inductive-formula-with-reduced}
    $a_{n+m} \cdot \beta^{\mathcal{A}_{n+m}^{m}} = \alpha^{n}\cdot\left(a_m + \frac{\gamma}{\alpha}\cdot\alpha^{m}\cdot\mathfrak{a}_{n+m}^{m}\right)$.
\end{corollary}
\begin{corollary}
Given $a\in\widehat{\mathbb{N}}_\beta$, $\mathfrak{A}$ is of $a$, we have $\beta \nmid \mathfrak{A}_n^m$.
\end{corollary}
\begin{corollary}~\label{corollary:congerunce-with-collatz-complement}
If $\mathcal{A}_k > 0$, then $a \equiv - \frac{\gamma}{\alpha}\cdot\mathfrak{a}_k = - \frac{\gamma}{\alpha^k}\cdot\mathfrak{A}_k \pmod{\beta^{\mathcal{A}_k}}$.
\end{corollary}
% \begin{proof} Trivial with \cref{corollary:inductive-formula-with-reduced}. \end{proof}
\begin{corollary}[Co-uniqueness]~\label{cor:co-uniqueness}
    If there are two numbers $a, b \in \widehat{\mathbb{N}}_{{\beta}}$ such that $\exists n \ge 0$, $\forall m \le n$, there is $\mathcal{A}_m = \mathcal{B}_m$, then $\beta^{\mathcal{A}_n} \mid (a - b)$ and $\alpha^{n} \mid (a_n - b_n)$.
\end{corollary}
% \begin{proof} Trivial with \cref{theorem:inductive-formula}. \end{proof}
\begin{proof}
    Since $\mathcal{A}_m = \mathcal{B}_m$, we have $\mathfrak{A}_{m} = \mathfrak{B}_{m}$.
    With the \cref{theorem:inductive-formula}, we have
    \[ (a_{m} - b_{m}) \cdot \beta^{\mathcal{A}_m} = (a - b) \cdot \alpha^{m}. \]
    Since $\gcd(\alpha, \beta) = 1$, then $\gcd(\beta^{\mathcal{A}_m}, \alpha^{m}) = 1$.
    So we have
    \[\beta^{\mathcal{A}_m} \mid (a - b), \quad \alpha^{m} \mid (a_m - b_m) \]
    holds for all $m\le n$.
\end{proof}
\begin{corollary}
    If $\forall m \ge 0$, there is $\mathcal{A}_m = \mathcal{B}_m$, then we have $a = b$.
\end{corollary}
\begin{corollary}
    If $a = t \cdot p^n + s$, that $t, s\in \mathbb{N}_{\ge1} - \beta\mathbb{N}_{\ge1}, p \in \mathbb{N}_{>1}$, then
    \[a_k \cdot \beta^{\mathcal{A}_k} - \gamma\cdot\mathfrak{A}_k = (t_k \cdot \beta^{\mathcal{T}_k} - \gamma\cdot\mathfrak{T}_k)\cdot p^{n} + (s_k \cdot \beta^{\mathcal{S}_k} - \gamma\cdot\mathfrak{S}_k).\]
\end{corollary}
\begin{corollary}
    If $a = t \cdot \beta^n + s$, that $t, s \in \widehat{\mathbb{N}}_{\beta;\ge1}$ and $n \in \mathbb{N}_{\ge1}$, then
    \begin{align*}
        \mathcal{A}_k &= \mathcal{S}_k, \\
        \mathfrak{A}_k &= \mathfrak{S}_k,\\
        a_k &= t \cdot \alpha^{k} \cdot \beta^{n-\mathcal{S}_k} + s_k
    \end{align*}
    holds during $S_k < n$.
\end{corollary}

It seems like we can introduce the Diophantine equation of $\mathfrak{A}_k$ now, but before that, we have to make sure the $\mathcal{A}_k$ will not finally increased by $0$, i.e.\ fall into those numbers cannot be reduced.
Though we have a clear evidence that $(\alpha, \beta, \gamma) = (3, 2, 1)$ guarantee that there exists at least one triple such that $\mathcal{A}_k \ge n$, we can't say all the triples required in \cref{definition:collatz-sequence} retain this property.
\begin{definition}
Given $m,n \in \mathbb{N}$, when $m \ge 1$, define
\[ [m\mid n] = \left\{a \in \mathbb{N}\ \left|\ \left(a\cdot\alpha^n + \gamma\cdot\sum_{i = 0}^{n-1}\alpha^i\right) \equiv 0 \pmod{\beta^{m}}\right.\right\}. \]
Specially, when $m = 0$, define $[0\mid n] = \mathbb{N}$.
\end{definition}
\begin{proposition}~\label{proposition:equivalence-of-classes}
    $a \in [m\mid n] \Longleftrightarrow a\equiv -\frac{\gamma}{\alpha}\cdot\sum_{i=0}^{n-1}\alpha^{-i}\pmod{\beta^{m}}$.
\end{proposition}
\begin{proposition}
    $ l \ge m \Longleftrightarrow [l\mid n] \subset [m\mid n]$.
\end{proposition}
\begin{proposition}~\label{proposition:avaliable-in-finity-step}%
    $\exists P \in \mathbb{N}_{\ge1}$, such that $[m \mid n + P ] = [m\mid n]$.
\end{proposition}
\begin{proof}
    Let $P = \ord_{\beta^{m}}(\alpha) \cdot \prod_{p\in\Pi(\mathbb{P} \mid \beta)}\raf_{p}(\alpha-1)$.

    First, since $\gcd(\alpha, \beta) = 1$, for those $p \in \mathbb{P}$ that $p \mid \beta$, we have
    \begin{align*}
        \fac_{p}(\alpha^P - 1) & \ge \fac_{p}\left(\prod_{p\in\Pi(\mathbb{P} \mid \beta)}\raf_{p}(\alpha-1)\right) + \fac_{p}(\alpha^{\ord_{\beta^{m}}(\alpha)} - 1) \\
                                 & \ge \fac_{p}(\alpha-1) + m \cdot \fac_{p}\beta,
    \end{align*}
    Then we have
    \begin{align*}
        \fac_{p}\left(\sum_{i=0}^{P - 1}\alpha^{-i}\right) & = \fac_{p}\left(\frac{\alpha(\alpha^{P} - 1)}{\alpha^{P}(\alpha - 1)}\right) \\
                &= \fac_{p}(\alpha^P - 1) - \fac_{p}(\alpha-1) \\
                &\ge m \cdot \fac_{p}\beta.
    \end{align*}
    That means $(\raf_{p_i}\beta)^{m} \left\vert\ \sum_{i=0}^{P - 1}\alpha^{-i} \right. $.
    Thus, we have $\beta^{m} \left\vert\ \sum_{i=0}^{P - 1}\alpha^{-i} \right. $.
    Here, $\alpha^{-1}$ is defined in terms of the modulus of $\beta^{m}$. \newline
    Now suppose $a \in [m \mid n + P]$.
    From \cref{proposition:equivalence-of-classes}, we have
    \begin{align*}
        a   &\equiv -\frac{\gamma}{\alpha}\cdot\sum_{i=0}^{P+n-1}\alpha^{-i}    &&\pmod{\beta^{m}} \\
            &\equiv -\frac{\gamma}{\alpha}\cdot\sum_{i=0}^{P - 1}\alpha^{-i} - \frac{\gamma}{\alpha}\cdot\sum_{i=P}^{P + n - 1}\alpha^{-i}    &&\pmod{\beta^{m}} \\
            &\equiv - \frac{\gamma}{\alpha}\cdot\alpha^{-P} \cdot\sum_{i=0}^{n - 1}\alpha^{-i}    &&\pmod{\beta^{m}} \\
            &\equiv -\frac{\gamma}{\alpha}\cdot\sum_{i=0}^{n - 1}\alpha^{-i}    &&\pmod{\beta^{m}}
    \end{align*}
    which means $a \in [ m \mid n + P] \Longleftrightarrow a \in [m \mid n]$.
    But this $P$ is not always the minimal period.
\end{proof}
\begin{definition}
    Denote $[m] = \bigcup_{n=0}^{\infty}[m\mid n]$ and $\bar{[m]} = \mathbb{N} - [m]$.
\end{definition}
\begin{proposition}
    $[0] = \mathbb{N}$ and $\bar{[0]} = \varnothing$.
\end{proposition}
\begin{proposition}
    $a \in \bar{[1]} \Longrightarrow \alpha \cdot a + \gamma \in \bar{[1]}$.
    On the contrary, if $a\equiv \gamma \pmod{\alpha}$, then $a \in \bar{[1]} \Longrightarrow \frac{a - \gamma}{\alpha} \in \bar{[1]}$.
\end{proposition}
\begin{theorem}~\label{theorem:requirements-of-avaliablity}
    $\lim_{k\to\infty} \mathcal{A}_k = \infty\Longleftrightarrow\forall k\in \mathbb{N}, a_k \in [1]$.
\end{theorem}
\begin{proof}
\begin{itemize}
    \item[$\Longrightarrow$)]
        If $\lim_{k\to\infty} \mathcal{A}_k = \infty$, then $\forall k \in \mathbb{N}$, $\exists n_{k} \in \mathbb{N}_{\ge1}$, $\val(a;k + n_{k}) > 0$,
        that means $a_{k} \in [\val(a;k + n_{k}) | n_{k}] \subset [1]$.
    \item[$\Longleftarrow$)] 
        If $\forall k \in \mathbb{N}$, there is $a_k \in [1]$, then $\exists n_{k} \in \mathbb{N}_{\ge1}$ and from \cref{proposition:avaliable-in-finity-step} we can suppose $n_k \le P$,
        such that $a_k \in [1 \mid n_{k}]$ for all $k \in \mathbb{N}$.
        Thus, we have $\mathcal{A}_k \ge \frac{k}{P} \to \infty\ (k \to \infty)$.
\end{itemize}
Sum up, we proved this theorem.
\end{proof}
\begin{corollary}
    If $\lim_{k\to\infty}\mathcal{A}_k = \infty$, then $\frac{\mathcal{A}_k}{k} \ge \frac{1}{P}$ when $k\to\infty$.
\end{corollary}
\begin{conjecture}
    $\exists c(\beta) \ge 0$, $\mu\{a \in\widehat{\mathbb{Z}}_{{\beta}} \mid \lim_{k\to\infty}\frac{\mathcal{A}_k}{k} \ge c \}  = 0$.
\end{conjecture}

\cref{theorem:requirements-of-avaliablity} present that, for general triplet arguments, not all initial numbers can be reduced infinitely many times.
But these propositions above provide a basic way to determine whether a number can be reduced, especially since we can easily filter out many initial numbers that cannot. We have
\begin{example}
    Suppose $(\alpha, \beta, \gamma) = (3, 2, 1)$.
    $\forall a \in \widehat{\mathbb{N}}_{2}$, $a\in[1]$, and $1 \le \frac{\mathcal{A}_k}{k} \le 2$ thus $\lim_{k\to\infty}\mathcal{A}_k = \infty$.
\end{example}

However, even if $A_{k}\to \infty$, we can't say that $a_k$ will eventually fall into a cycle.
For example, when $\alpha = 11$, $\beta = 2$ and $\gamma = 1$, for any number, the sequence may grow exponentially as $k$ increases, while $\alpha = 3$ will decrease to 1, at least for the numbers we observe.
The reason for this is that the distribution of reduce always only relates to $\beta$, and its change with $\alpha$ is not significant.
But the growth of sequence numbers accelerates as $\alpha$ increases. $11$ is too large, $\beta=2$ and $\gamma=1$ are dwarfed with it.

Therefore, we can only find these cycle numbers by letting $a_{t} = a$.
According to \cref{theorem:inductive-formula}, that is
\[ a = \frac{\gamma \cdot \mathfrak{A}_{t}}{\beta^{\mathcal{A}_{t}} - \alpha^{t}}. \]
For this equation, there probably are some integer solutions from \cref{proposition:trivial-periodic-sequence}, but not all integer solution of it follows that proposition.
For example, $a=13$ has a 4-tempus cycle under $\alpha=3$, $\beta=5$, $\gamma=1$, that $\val(13;1) = 1$, $\val(13;2) = 2$, $\val(13;3) = 0$ and $\val(13;4) = 0$.
Obviously, $\val(13;i)$ are not same.
Therefore, we have to explore the properties of this fractional expression independently.

% !TEX root = main.tex
\section{Extend the Field of Values for the Collatz Sequence}\label{sec:extend-the-field-of-values-for-the-collatz-sequence}

A better thought is extending our view from $\mathbb{N}$ to $\mathbb{Q}$.
Keeping $\alpha, \beta, \gamma \in \mathbb{N}$, instead of restricting the initial number $a \in \mathbb{N}$, we can let $a$ be any number of $\mathbb{Q}-\beta\mathbb{Q}$.
In this situation, the \cref{theorem:inductive-formula} still holds, so do those corollaries of the theorem.
More than that, we can get a Collatz sequence by given the reduces.

\begin{theorem}~\label{theorem:any-cycle}
    Given $\{v_{i}\}_{i=1}^{n}$ that $n < \infty$, if $\forall i>j\ge1, \beta \nmid \mathfrak{A}_i^{j}$, then there exists only one Collatz sequence in $\mathbb{Q}$ is a cycle that $\val_{\beta}(a;k) = v_{k\fmod n}$.
\end{theorem}
\begin{proof}%[Proof of \cref{theorem:any-cycle}]
    Let's just using the notation $\mathcal{A}_k$ and $\mathfrak{A}_k$ as if $\val_{\beta}(a;k) = v_{k\fmod n}$.\newline
    Then $\forall k = 1, 2, \cdots, n$, we have $\mathcal{A}_{k + n}^{k} = \mathcal{A}_n$.
    Let
    \[ a_k = \frac{\gamma\cdot\mathfrak{A}_{n+k}^{k}}{\beta^{\mathcal{A}_n} - \alpha^n} \in \mathbb{Q}. \]
    Obviously, we have $\beta \nmid a_{k}$.
    And we have
    \begin{align*}
        \alpha\cdot a_{k} + \gamma &= \frac{\gamma\cdot\alpha\cdot\mathfrak{A}_{n+k}^{k}}{\beta^{\mathcal{A}_n} - \alpha^n} + \frac{\gamma\cdot(\beta^{\mathcal{A}_n} - \alpha^n)}{\beta^{\mathcal{A}_n} - \alpha^n} \\
        &= \frac{\gamma\cdot(\alpha\cdot\mathfrak{A}_{n+k}^{k}+\beta^{\mathcal{A}_{n+k}^{k}}) - \gamma\cdot\alpha^n}{\beta^{\mathcal{A}_n} - \alpha^n} \\
        &= \frac{\gamma\cdot(\mathfrak{A}_{n+k+1}^{k} - \alpha^{n})}{\beta^{\mathcal{A}_n} - \alpha^n} \\
        &= \frac{\gamma\cdot\beta^{\mathcal{A}_{k+1}^{k}}\cdot\mathfrak{A}_{n+k+1}^{k+1}}{\beta^{\mathcal{A}_n} - \alpha^n} \\
        &= \frac{\gamma\cdot\mathfrak{A}_{n+k+1}^{k+1}}{\beta^{\mathcal{A}_n} - \alpha^n}\cdot \beta^{\mathcal{A}_{k+1}^{k}} \\
        &= a_{k+1}\cdot \beta^{\Delta_{k+1}}.
    \end{align*}
    So we just let $a = a_{n}$, then the Collatz sequence of $a$ is what we want.
    And it is unique.
\end{proof}
In fact, we can see that the rotation of $v_{k}$ corresponds with the rotation of $a_k$.
That implies some symmetry of the Collatz sequence.
But not all Collatz reduce is a pure cycle.
In fact, most of  what may eventually enter a cycle have a chaotic beginning.
Even there may exist some Collatz sequence will not fall into any cycle.
For these two kinds of Collatz reduces, we have
\begin{theorem}~\label{theorem:any-presequence}
    Given $\{v_{i}\}_{i=1}^{n}$ that $n < \infty$, and any number $b$.
    If $\forall i>j\ge1, \beta \nmid \mathfrak{A}_i^{j}$, then there exists only one Collatz sequence in $\mathbb{Q}[b]$ that $\val_{\beta}(a;i) = v_{i}$ and $a_n = b$.
\end{theorem}
\begin{proof}
    Also using the notation $\mathcal{A}_k$, $\mathfrak{A}_k$ and $\mathfrak{a}_k$ as if $\val_{\beta}(a;i) = v_{i}$.
    Let
    \[ a_k = \frac{b \cdot \beta^{\mathcal{A}_n^{k}} - \gamma\cdot\mathfrak{A}_n^{k}}{\alpha^{n-k}} = b \cdot \frac{\beta^{\mathcal{A}_n^{k}}}{\alpha^{n-k}}  - \frac{\gamma}{\alpha}\cdot\alpha^{k}\cdot\mathfrak{a}_n^{k} \in \mathbb{Q} \]
    where $k = 0, 1, 2, \cdots, n$.
    Similarly, since $\beta \nmid \mathfrak{A}_n^k$, then $\beta \nmid \mathfrak{a}_n^{k}$, then $\beta \nmid a_k$.
    And we have
    \begin{align*}
        \alpha \cdot a_k + \gamma &= b \cdot \frac{\beta^{\mathcal{A}_n^{k}}}{\alpha^{n-k-1}}  - \gamma\cdot\alpha^{k}\cdot\mathfrak{a}_n^{k} + \gamma \\
        % &= b \cdot \frac{\beta^{\mathcal{A}_n^{k}}}{\alpha^{n-k-1}}  - \gamma\cdot\alpha^{k}\cdot\left(\mathfrak{a}_n^{k} - \frac{1}{\alpha^{k}}\right) \\
        &= b \cdot \frac{\beta^{\mathcal{A}_n^{k}}}{\alpha^{n-k-1}}  - \gamma\cdot\alpha^{k}\cdot\left(\mathfrak{a}_n^{k} - \mathfrak{a}_{k+1}^{k}\right) \\
        &= \left(b \cdot \frac{\beta^{\mathcal{A}_n^{k+1}}}{\alpha^{n-(k+1)}}  - \frac{\gamma}{\alpha}\cdot\alpha^{k+1}\cdot\mathfrak{a}_{n}^{k+1}\right) \cdot \beta^{\mathcal{A}_{k+1}^{k}} \\
        &= a_{k+1} \cdot \beta^{\Delta_{k+1}}
    \end{align*}
    to make sure it is a Collatz sequence.
    Finally, we have $a_n = b$.
    Also, it is unique.
\end{proof}
\begin{corollary}
    Given $a = \frac{   p}{q}\in \mathbb{Q}$. If $\exists s\in \mathbb{N}$, such that $q = \alpha^{s}$, then $a_{s} \in \mathbb{N}$ by natual reduce sequence.
\end{corollary}

\cref{theorem:any-cycle} and \cref{theorem:any-presequence} present that we can come up with a lot of Collatz sequence that will fall into cycles by given its reduce sequence instead of startswith an integer.
However, not all given reduce sequence can derive a Collatz sequence because there may be some $\beta\mid \mathfrak{A}_{n+k}^{k}$.
At that time, we have $\beta \left\vert \frac{\gamma\cdot\mathfrak{A}_{n+k}^{k}}{\beta^{\mathcal{A}_n} - \alpha^n} \right.$.
That's why we need $\beta \nmid \mathfrak{A}_i^j$.
Moreover, from the two proofs above, we can also draw another interesting corollary:
\begin{theorem}
    If Collatz sequence of $a$ will fall into a cycle, then $a \in \mathbb{Q}$.
\end{theorem}
However, the converse of this conclusion does not hold.
For example, $a=3$ under $\alpha = 11$, $\beta=2$ and $\gamma=1$ may never fall into a cycle.
For these won't fall in to cycles, there is also a conjecture:
\begin{conjecture}
    Given $\{v_{i}\}_{i=1}^{\infty}$ has no cycle, then there is a Collatz sequence that $\val_{\beta}(a;i) = v_{i}$, it has a subsequence $\{a_{k_i}\}_{i=0}^{\infty}$ converges to an $\beta$-adic irrational number or $0$.
\end{conjecture}

Anyway, we now finally extended the problem from integers to the rational numbers, and at the same time, the exploration of Collatz sequences has turned to the exploration of Collatz reduces.
This implies that the Collatz reduces is the more fundamental aspect of this process.

% !TEX root = main.tex
\section{Back Trace the Collatz Sequence}\label{sec:back-trace-the-collatz-sequence}
\cref{theorem:any-presequence} tell us we can also trace back the Collatz sequence in $\mathbb{Q}$.
But we can only go back under a finite reduces given.
For an infinite tracing back, we need to redefine the summation to give it directionality.
\begin{definition}
    Let's modified the definition of $\sum$ operator. The sum $S$ of $s$ from $m$ to $n$ is defined as
    \[ S = \left\{\begin{aligned}
    &\sum_{i=m}^{n-1}s_i, &&n > m, \\
    &0, &&n = m, \\
    -&\sum_{i=n}^{m-1}s_i, &&n < m.
\end{aligned}\right. \]
\end{definition}
With this we have
\begin{proposition}
    $\mathcal{A}_{n}^{m} = -\mathcal{A}_{m}^{n}$.
\end{proposition}
\begin{proposition}
    $\alpha^{n-m} \cdot \mathfrak{A}_{m}^{n} + \beta^{\mathcal{A}_m^n} \cdot \mathfrak{A}_{n}^{m} = 0$ and $\mathfrak{a}_{m}^{n} + \beta^{\mathcal{A}_m^n} \cdot \mathfrak{a}_{n}^{m} = 0$.
\end{proposition}
It is easy to verify that those propositions of Collatz accumulation and complement still holds.
And with these definitions, we have
\begin{theorem}~\label{theorem:final-inductive-formula}
Given $a \in \widehat{\mathbb{Q}}_{{\beta}}$ and $\{v_{i}\}_{i=-\infty}^{\infty} \subset \mathbb{N}$ that extends from the Collatz reduces of $a$ and $\forall i, j \in \mathbb{Z}$, there is $\beta \nmid \mathfrak{A}_{i}^{j}$.
Let $a_0 = a$, then $\forall m, n \in \mathbb{Z}$, there still is
\[ a_n \cdot \beta^{\mathcal{A}_{n}^{m}} = a_m \cdot\alpha^{n-m} + \gamma\cdot\mathfrak{A}_{n}^{m}. \]
\end{theorem}

\begin{corollary}
    $a_n \cdot \beta^{\mathcal{A}_{n}^{m}} = \alpha^{n-m} \cdot\left(a_m + \frac{\gamma}{\alpha}\cdot\alpha^{m}\cdot\mathfrak{a}_{n}^{m}\right)$.
\end{corollary}
\begin{corollary}
    If $\forall i$ that $k \ge i > 0$, there is $v_{-i} = t$, denote $\sigma = \frac{\gamma}{\alpha-\beta^{t}}$, then we have
    $(a_{-k} + \sigma) \cdot \alpha^k  = \beta^{kt} \cdot (a_0 + \sigma)$.
\end{corollary}
Equation $(a_{-k} + \sigma) \cdot \alpha^k  = \beta^{kt} \cdot (a_0 + \sigma)$ has only one form of solution:
\[\left\{\begin{aligned}
    a_{0} &= t \cdot \alpha^k - \sigma, \\
    a_{-k} &= t \cdot \beta^{k\Delta} - \sigma,
\end{aligned}\right.\]
where $t \in \mathbb{Q}$.
To ensure both $a_{-k}$ and $a_0$ are integers, we require
\[ t \in \{s\cdot\sigma+n \mid n \in \mathbb{N}, s \equiv \alpha^{-k} \equiv \beta^{-kt} \pmod{\alpha - \beta^{t}}\}. \]
Unfortunately, in this case, not all solutions of the equation guarantee that $\beta \nmid a_0$.

% !TEX root = main.tex
\section{Varies in Source Value with Respect to the Reduce Sequence}\label{sec:varies-in-source-value-with-respect-to-the-reduce-sequence}

\begin{definition}
From \cref{theorem:any-presequence}, with the same arguments $\alpha$, $\beta$ and $\gamma$, and target value $b\in \mathbb{Q} - \beta\mathbb{Q}$, for any finite vector $\bm{v} = (v_{1}, v_{2}, \cdots, v_{n}) \in \mathbb{N}^{n}$, we have a function $f_n : \mathbb{N}^{n} \to \mathbb{Q}$, that defined as
\[ f_{n}(\bm{v};b) = b \cdot \frac{\beta^{\mathcal{V}_{n}}}{\alpha^{n}} - \frac{\gamma}{\alpha}\cdot \mathfrak{v}_{n} \in \mathbb{Q} \]
is called the source value from $b$ by $\bm{v}$. In the absence of ambiguity, we denote $f_{n}(\bm{v}) = f_{n}(\bm{v};b)$. Given $s \in \mathbb{Z}$, there is a mapping $F_{s; k}: \mathbb{N}^{n} \to \mathbb{N}^{n}$ that
\[ F_{s; k}(\bm{v}) = \bm{v} + s\bm{e}_k = (v_{1}, v_{2}, \cdots,v_{k} + s, \cdots, v_{n}), \]
where $\bm{e}_k$ is the $k$-th unit vector.
And we also sometimes denote $F_{s; k}(\mathfrak{V})$ and $F_{s; k}(\mathfrak{v})$ to present those Collatz complements and reduced Collatz complements that from $F_{s; k}(\bm{v})$.
\end{definition}

And to describe the iterative properties of Collatz sequences, we denote $a \overset{k}{\rightsquigarrow} b$ if there exists $k$ such that $a_k = b$.
If it $k$ doesn't matter, we denote $a \rightsquigarrow b$.
Immediately, we have
\begin{proposition}
    $f_{n}(\bm{v}) \overset{n}{\rightsquigarrow} b$.
\end{proposition}

\begin{proposition}
    If $\mathfrak{v}_{n}$ follow the requirements of \cref{theorem:any-presequence}, then we will have $\beta \nmid f_{n}(\bm{v})$.
\end{proposition}
\begin{proposition}
    $f_{n}(\bm{v}) = f_{n}(\bm{u}) \Longleftrightarrow \bm{v} = \bm{u}$.
\end{proposition}
\begin{proposition}
    $f_{n}(\bm{v}) = \frac{\beta^{v_{1}}}{\alpha}\left(\frac{\beta^{v_{2}}}{\alpha}\left(\cdots \left(\frac{\beta^{v_{n}}}{\alpha}\cdot b - \frac{\gamma}{\alpha}\right)\cdots\right)- \frac{\gamma}{\alpha}\right) - \frac{\gamma}{\alpha}$.
\end{proposition}
\begin{proposition}
    $F_{t;l}(F_{s;k}(\bm{v})) = F_{s;k}(F_{t;l}(\bm{v}))$.
\end{proposition}
\begin{proposition}~\label{prop:single-vary-rule}
$f_{n}(F_{s;k}(\bm{v})) = \beta^{s}\cdot f_{n}(\bm{v}) + \frac{\gamma}{\alpha}\cdot\mathfrak{v}_{k}\cdot(\beta^{s} - 1)$.
\end{proposition}
\begin{proposition}
    Suppose $m\le n$, then
    \begin{align*}
        F_{s;k}(\mathcal{V})_{n}^{m} &= \left\{\begin{aligned}
        &\mathcal{V}_{n}^{m} + s, & m\le k < n, \\
        &\mathcal{V}_{n}^{m}, & \text{otherwise},
    \end{aligned}\right. \\
    F_{s;k}(\mathfrak{V})_{n}^{m} &= \left\{\begin{aligned}
        &\mathfrak{V}_{k}^{m} + \beta^{s}\cdot \mathfrak{V}_{n}^{k}, & m\le k \le n, \\
        &\mathfrak{V}_{n}^{m}, & \text{otherwise},
    \end{aligned}\right. \\
    F_{s;k}(\mathfrak{v})_{n}^{m} &= \left\{\begin{aligned}
        &\mathfrak{v}_{k}^{m}\cdot\alpha^{k-n} + \beta^{s}\cdot \mathfrak{v}_{n}^{k}, & m\le k \le n, \\
        &\mathfrak{v}_{n}^{m}, & \text{otherwise}.
    \end{aligned}\right.
    \end{align*}
\end{proposition}
\begin{proposition}
$\frac{\gamma}{\alpha}\cdot\mathfrak{v}_{k} = \frac{f_{n}(F_{s;k}(\bm{v})) - \beta^{s}\cdot f_{n}(\bm{v})}{\beta^{s} - 1}$.
\end{proposition}
\begin{proposition}
    If $f_{n}(\bm{v}) = \frac{q}{\alpha^{m}}$, then $f_{n}(F_{s;k}(\bm{v})) \in \mathbb{N}$ only when $s$ holds
    \[ \beta^{s}\cdot q\cdot \alpha^{k} + (\beta^{s} - 1) \cdot \mathfrak{V}_{k}\cdot \alpha^{m} \equiv 0. \pmod{\alpha^{m+k}} \]
\end{proposition}
\begin{theorem}~\label{thm:vary-rule}
Given $\Delta\bm{v} \in \mathbb{Z}^{n}$, view it as a $n$-length Collatz sequence. Then we have
\[ f_{n}(\bm{v} + \Delta\bm{v}) = \beta^{\Delta\mathcal{V}_{n}}\cdot f_{n}(\bm{v}) + \frac{\gamma}{\alpha^{n}} \cdot \nabla\cdot\mathfrak{V}_{n} = \beta^{\Delta\mathcal{V}_{n}}\cdot f_{n}(\bm{v}) + \frac{\gamma}{\alpha} \cdot \nabla\cdot\mathfrak{v}_{n}. \]
Where
\begin{align*}
    \nabla\cdot\mathfrak{V}_{n} &= \sum_{k=0}^{n-1}\beta^{\mathcal{V}_{k} + \Delta\mathcal{V}_{k}}\cdot\alpha^{n-k-1}\cdot(\beta^{\Delta\mathcal{V}_{n}^{k}} - 1),\\
    \nabla\cdot\mathfrak{v}_{n} &= \sum_{k=0}^{n-1}\frac{\beta^{\mathcal{V}_{k} + \Delta\mathcal{V}_{k}}}{\alpha^{k}}\cdot(\beta^{\Delta\mathcal{V}_{n}^{k}} - 1).
\end{align*}
\end{theorem}
\begin{proof}
    Let $\bm{u} = \bm{v} + \Delta\bm{v}$. First we have
    \begin{align*}
        \mathfrak{U}_{n} &= \beta^{\Delta\mathcal{V}_n}\cdot\mathfrak{V}_n - \nabla\cdot\mathfrak{V}_{n},\\
        \mathfrak{u}_{n} &= \beta^{\Delta\mathcal{V}_n}\cdot\mathfrak{v}_n - \nabla\cdot\mathfrak{v}_{n}.
    \end{align*}
    Notice that
    \[f_{n}(\bm{u}) = b \cdot \frac{\beta^{\mathcal{U}_{n}}}{\alpha^{n}} - \frac{\gamma}{\alpha}\cdot \mathfrak{u}_{n} = \beta^{\Delta\mathcal{V}_n} \cdot b\cdot \frac{\beta^{\mathcal{V}_{n}}}{\alpha^{n}} - \frac{\gamma}{\alpha}\cdot \mathfrak{u}_{n},\]
    and
    \[ \mathfrak{u}_{n} = \sum_{k=0}^{n-1}\frac{\beta^{\mathcal{U}_k}}{\alpha^{k}} = \sum_{k=0}^{n-1}\frac{\beta^{\mathcal{V}_k + \Delta\mathcal{V}_k}}{\alpha^{k}}. \]
    Then
    \begin{align*}
        & f_{n}(\bm{u}) - \beta^{\Delta\mathcal{V}_{n}}\cdot f_{n}(\bm{v}) \\
        =& \frac{\gamma}{\alpha}\cdot\left(\beta^{\Delta\mathcal{V}_{n}}\cdot\mathfrak{v}_{n} - \mathfrak{u}_{n} \right) \\
        =&\frac{\gamma}{\alpha^{n}}\cdot\left(\beta^{\Delta\mathcal{V}_{n}}\cdot\sum_{k=0}^{n-1}\beta^{\mathcal{V}_{k}}\cdot\alpha^{n-k-1} - \sum_{k=0}^{n-1}\beta^{\mathcal{V}_{k}+\Delta\mathcal{V}_{k}}\cdot\alpha^{n-k-1} \right) \\
        =&\frac{\gamma}{\alpha^{n}}\cdot\sum_{k=0}^{n-1}\beta^{\mathcal{V}_{k}}\cdot(\beta^{\Delta\mathcal{V}_{n}} - \beta^{\Delta\mathcal{V}_{k}})\cdot\alpha^{n-k-1} \\
        =&\frac{\gamma}{\alpha^{n}}\cdot\sum_{k=0}^{n-1}\beta^{\mathcal{V}_{k}}\cdot\alpha^{n-k-1}\cdot\beta^{\Delta\mathcal{V}_{k}}\cdot(\beta^{\Delta\mathcal{V}_{n}^{k}} - 1) \\
        =&\frac{\gamma}{\alpha^{n}}\cdot\sum_{k=0}^{n-1}\beta^{\mathcal{V}_{k} + \Delta\mathcal{V}_{k}}\cdot\alpha^{n-k-1}\cdot(\beta^{\Delta\mathcal{V}_{n}^{k}} - 1)\\
        =&\frac{\gamma}{\alpha^{n}}\cdot \nabla\cdot\mathfrak{U}_{n}
        \\
        =&\frac{\gamma}{\alpha}\cdot \nabla\cdot\mathfrak{u}_{n}.
    \end{align*}
    Finally, this proposition has been proved.
\end{proof}
Though $\Delta\bm{v} \in \mathbb{Z}^{n}$, we still require $\bm{v} + \Delta\bm{v} \in\mathbb{N}^{n}$.
If there is a vector $\Delta\bm{v}$ with only one non-zero component, i.e. $\bm{v} + \Delta\bm{v}$ becomes $F_{s;k}(\bm{v})$, then the analysis becomes easier.
\begin{theorem}[The Periodic Rule of Reduce Sequence]~\label{thm:the-periodic-rule-of-reduce-sequence}
    If $f_{n}(\bm{v}) \in \mathbb{N}$, then
    \[f_{n}(F_{s;k}(\bm{v})) \in \mathbb{N}\Longleftrightarrow \left.\frac{\alpha^{k}}{\raf_{\alpha}(\gamma\cdot2^{\delta})}\cdot\Psi_{\alpha}(\beta)\right\vert s, \]
    where the auxiliary operator $\Psi_{\alpha}(\beta)$ is defined as \[
        \Psi_{\alpha}(\beta) = \underset{p\mid\alpha}{\lcm}\{\ord_{p}(\beta)\}\cdot\prod_{p\in\mathbb{P}}^{p \mid \alpha}\frac{1}{\raf_{p}(\beta^{\ord_{p}(\beta)} - 1)},
    \]
    and $\delta = \fac_2(\beta + 1)$ if $2 \mid \alpha$ else $0$.
\end{theorem}
\begin{proof}
    From \cref{prop:single-vary-rule}, we have
    \[ f_{n}(F_{s;k}(\bm{v})) = \beta^{s}\cdot f_{n}(\bm{v}) + \frac{\gamma}{\alpha}\cdot\mathfrak{v}_{k}\cdot(\beta^{s} - 1). \]
    Since $\beta^{s}\cdot f_{n}(\bm{v}) \in \mathbb{N}$, then $\frac{\gamma}{\alpha}\cdot\mathfrak{v}_{k}\cdot(\beta^{s} - 1) \in \mathbb{N}$, that equivalents
    \[\alpha^{k} \mid \gamma\cdot\mathfrak{V}_{k}\cdot(\beta^{s} - 1).\]
    Since $\gcd(\alpha, \mathfrak{V}_{k}) = 1$, the divisibility condition simplifies to
    \[\alpha^{k} \mid \gamma\cdot(\beta^{s} - 1).\]
    By the Fundamental Theorem of Arithmetic, $\alpha = \prod_{p \in \mathbb{P}}^{p \mid \alpha} p^{\fac_{p}(\alpha)}$. Thus, the global condition is equivalent to a set of local conditions for each prime $p \mid \alpha$ we have
    \[ p^{k\cdot\fac_{p}(\alpha)} \mid \gamma\cdot(\beta^{s} - 1). \]
    Note that $p \mid (\beta^{\ord_{p}(\beta)} - 1)$, apply the Lifting The Exponent (LTE) Lemma, consequently, the inequality becomes
    \[ k\cdot\fac_{p}(\alpha) \le \fac_{p}(\gamma) + \fac_{p}\left(\frac{s}{\ord_{p}(\beta)}\right) + \fac_{p}(\beta^{\ord_{p}(\beta)} - 1) + \delta_2(\beta). \]
    Here $\delta_2(p) = \fac_2(\beta + 1)$ if $p = 2$ else $0$.
    Rearranging for $\fac_{p}\left(\frac{s}{\ord_{p}(\beta)}\right)$, we obtain
    \[ \fac_{p}\left(\frac{s}{\ord_{p}(\beta)}\right) \ge \fac_{p}(\alpha^{k}) - \fac_{p}(\gamma) - \fac_{p}(\beta^{\ord_{p}(\beta)} - 1) - \delta_2(\beta). \]
    Isolating the valuation of $s$, we obtain
    \[ \fac_{p}(s) \ge \fac_{p}\left(\frac{\alpha^{k}\cdot\ord_{p}(\beta)}{\gamma\cdot(\beta^{\ord_{p}(\beta)} - 1)\cdot2^{\delta}}\right), \]
    which equivalents
    \[ \left.\raf_{p}\left(\frac{\alpha^{k}\cdot\ord_{p}(\beta)}{\gamma\cdot(\beta^{\ord_{p}(\beta)} - 1)\cdot2^{\delta}}\right) \right\vert s. \]
    From \cref{prop:raf-completely-multiplicative}, we have
    \[ \raf_{p}\left(\frac{\alpha^{k}\cdot\ord_{p}(\beta)}{\gamma\cdot(\beta^{\ord_{p}(\beta)} - 1)\cdot2^{\delta}}\right) = \frac{p^{k\cdot\fac_{p}(\alpha)}}{\raf_{p}(\gamma\cdot2^{\delta})}\cdot\frac{\raf_{p}\ord_{p}(\beta)}{\raf_{p}(\beta^{\ord_{p}(\beta)} - 1)}, \]
    which implies the divisibility
    \[ \left. \frac{p^{k\cdot\fac_{p}(\alpha)}}{\raf_{p}(\gamma\cdot2^{\delta})}\cdot\frac{\raf_{p}\ord_{p}(b)}{\raf_{p}(\beta^{\ord_{p}(\beta)} - 1)} = \frac{p^{k\cdot\fac_{p}(\alpha)}}{\raf_{p}(\gamma\cdot2^{\delta})}\cdot\frac{1}{\raf_{p}(\beta^{\ord_{p}(\beta)} - 1)} \right\vert s. \]
    Besides, we've assumed $\ord_{p}(\beta) \mid s$. Aggregate these all $p$ conditions, it follows
    \[ \underset{p\mid \alpha}{\lcm}\{ \ord_{p}(\beta) \} \mid s. \]
    Sum up, we have
    \[ \left.\frac{\alpha^{k}}{\raf_{\alpha}(\gamma\cdot2^{\delta})}\cdot\Psi_{\alpha}(\beta)\right\vert s. \]
\end{proof}
\cref{thm:the-periodic-rule-of-reduce-sequence} represents the changes $s$ that preserve integers, which do not depend on the original reduced sequence $\bm{v}$, but only on the position $k$.
Therefore, we can classify these reduced sequences into different equivalence classes.

Furthermore, this conclusion can be easily generalized to the study of the periodic properties of linear or fractional dynamical systems in other algebraic fields, simply by replacing the prime factors with ideals and the LCM with the least common ideal, while the form remains essentially unchanged.
However, such generalizations are beyond the scope of this paper, which will continue to focus on Collatz sequences.

Back to the \cref{thm:vary-rule}.
This proposition offers a very important insight.
That is, if $f_{n}(\bm{v})\overset{n}{\rightsquigarrow} b$ and $f_{n}(\bm{v} + \Delta\bm{v}) \overset{n}{\rightsquigarrow} b$, then starting from a certain bit $m$ and $m + \Delta\mathcal{V}_n$, the $\beta$-carrying terms of these two numbers satisfy a certain relationship. Mathematically, that's
\begin{theorem}~\label{thm:high-order-bits-sync}
    Suppose $f_{n}(\bm{v}) = \sum_{k=s}^{\infty}x_k\beta^{k}$ and $f_{n}(\bm{v} + \Delta\bm{v}) = \sum_{k=s}^{\infty}x'_k\beta^{k}$, $s\in\mathbb{Z}$. If $\alpha^{n} \mid \gamma\cdot\nabla\cdot\mathfrak{V}_{n}$, then $\exists m \in \mathbb{Z}: m \ge s$, such that
    \[ \sum_{k=0}^{\infty}x'_{k + m}\beta^{k} = \sum_{k=0}^{\infty}x_{k + m - \Delta\mathcal{V}_n}\beta^{k} + \delta, \]
    where the $\delta\in\{0, 1\}$. Vice versa.
\end{theorem}
\begin{proof}
    When $\nabla\cdot\mathfrak{V}_{n} = 0$, it's trivial. Consider $\nabla\cdot\mathfrak{V}_{n} \ne 0$.\newline
    Since $\alpha^{n} \mid \gamma\cdot\nabla\cdot\mathfrak{V}_{n}$, then $\beta^{\vert\Delta\mathcal{V}_n\vert}\cdot\frac{\gamma}{\alpha}\cdot\nabla\cdot\mathfrak{v}_{n} \in \mathbb{N}$, that means $\deg_{\beta}(\frac{\gamma}{\alpha}\cdot\nabla\cdot\mathfrak{v}_{n}<)\infty$. Now let $m = \max\{\deg_{\beta}(\frac{\gamma}{\alpha}\cdot\nabla\cdot\mathfrak{v}_{n}), \Delta\mathcal{V}_n\} + s + 1$, then
    \[ \beta^{m - 1}\le\beta^{\Delta\mathcal{V}_n}\sum_{k=s}^{m - \Delta\mathcal{V}_n - 1}x_k\beta^{k} + \frac{\gamma}{\alpha}\cdot\nabla\cdot\mathfrak{v}_{n} \le \beta^{m}. \]
    And from \cref{thm:vary-rule}, we have
    \[ \sum_{k=s}^{\infty}x'_k\beta^{k} = \beta^{\Delta\mathcal{V}_n}\sum_{k=s}^{m - \Delta\mathcal{V}_n - 1}x_k\beta^{k} + \frac{\gamma}{\alpha}\cdot\nabla\cdot\mathfrak{v}_{n} + \beta^{\Delta\mathcal{V}_n}\sum_{k=m}^{\infty}x_{k}\beta^{k}. \]
    Let's right shift both the series by $m$ then we have
    \[ \sum_{k=m}^{\infty}x'_k\beta^{k - m} = \sum_{k=m - \Delta\mathcal{V}_n}^{\infty}x_{k}\beta^{k-m + \Delta\mathcal{V}_{n}} + \delta. \]
    Adjusting the subscript yields
    \[ \sum_{k=0}^{\infty}x'_{k + m}\beta^{k} = \sum_{k=0}^{\infty}x_{k + m - \Delta\mathcal{V}_n}\beta^{k} + \delta. \]
    On the other hand, if there is $m$ make the series equation holds, then $\forall k \ge m$ the series equation also holds. So suppose $m \ge \Delta\mathcal{V}_n + s + 1$. We can see from the previous derivation that
    \[\frac{\gamma}{\alpha}\cdot\nabla\cdot\mathfrak{v}_{n} = \sum_{k=s}^{m-1}x'_k\beta^{k} - \beta^{\Delta\mathcal{V}_n}\sum_{k=s}^{m - \Delta\mathcal{V}_n - 1}x_k\beta^{k} + \delta\cdot\beta^{m}, \]
    thus $\deg_{\beta}(\frac{\gamma}{\alpha}\cdot\nabla\cdot\mathfrak{v}_{n}) \le m \le\infty$. Notice that $\gcd(\alpha, \beta) = 1$, thus $\alpha^{n}\mid \gamma\cdot\nabla\cdot\mathfrak{V}_{n}$.
\end{proof}
Notice that we don't require the $f_{n}(\bm{v})$ and $f_{n}(\bm{v} + \Delta\bm{v})$ in $\mathbb{N}$, and we don't require $m \ge 0$. that means they can be any $\beta$-adic integers or $\beta$-adic finite decimals.
\begin{remark}
    If $f_{n}(\bm{v}) \not\in \mathbb{Z}_{<0}$, then $\exists m > \max\{\deg_{\beta}(\frac{\gamma}{\alpha}\cdot\nabla\cdot\mathfrak{v}_{n}), \Delta\mathcal{V}_n\} + s + 1$, such that \cref{thm:high-order-bits-sync} holds and $\delta = 0$. But if $f_{n}(\bm{v}) \in \mathbb{Z}_{<0}$ and there is
    \[ \frac{\gamma}{\alpha}\cdot\nabla\cdot\mathfrak{v}_{n} \ge \left\vert \beta^{\Delta\mathcal{V}_n}\cdot f_{n}(\bm{v}) \right\vert, \]
    then $\forall m$ large enough, there is only $\delta = 1$, i.e. $f_{n}(\bm{v} + \Delta\bm{v}) \in \mathbb{N}$.
\end{remark}
Actually, we can derive this theorem from \cref{theorem:any-presequence}. Since $\gcd(\alpha, \beta) = 1$, $\alpha^{-n}$ is a rational number, $\alpha^{n} \mid \gamma\cdot\nabla\cdot\mathfrak{V}_{n}$ means $\frac{\gamma}{\alpha}\cdot\nabla\cdot\mathfrak{v}_{n} \in \mathbb{N}$ and will not change the fraction part of the mixed fraction expression of $\beta^{\Delta\mathcal{V}_{n}}\cdot f_{n}(\bm{v}) + \frac{\gamma}{\alpha} \cdot \nabla\cdot\mathfrak{v}_{n}$. In $\beta$-adic integer $\mathbb{Z}_{\beta}$, it means the highter bits will not changed. And the $\beta^{\Delta\mathcal{V}_{n}}$ determinate the shift count, it also will not change the highter bits.

Even though, we can't determinate integer source values only depends on $\nabla\cdot\mathfrak{V}_{n}\in\mathbb{N}$, because there are evidences present that even if $\nabla\cdot\mathfrak{V}_{n}\not\in\mathbb{N}$, we can derive an integer source value.

Besides these varies in source value, just from the definition of source value, we have a two kinds of important sets.
\begin{definition}
Given a basic target value set $B$ denote that
\begin{align*}
    \mathbb{V}_n(B) &= \bigcup_{b\in B}\left\{ \bm{v} \in \mathbb{N}^{n} \mid f_{n}(\bm{v};b) \in \mathbb{N} \right\},\\
    \mathbb{F}_n(B) &= \bigcup_{b\in B}\left\{ f_n(\bm{v};b) \mid \bm{v} \in \mathbb{A}_n(b) \right\}.
\end{align*}
Specially, we define $\mathbb{V}_0(B) = \varnothing$ and $\mathbb{F}_0(B) = B$. When $B = \{b\}$, that's $B$ contains only one number, we also denote as $\mathbb{V}_n(b)$ and $\mathbb{F}_n(b)$.
\end{definition}
\begin{proposition}
If $b$ in a cycle which period is $t$, then $\forall k\ge 0$, $b \in \mathbb{F}_{kt}(b)$, thus $\mathbb{F}_{kt}(b)\subset \mathbb{F}_{(k+1)\cdot t}(b)$.
\end{proposition}
Actually we can derive $\mathbb{A}_{n+1}(b)$ from $\mathbb{A}_{n}(b)$. But since we are calculating the source value, the positive direction of the vector's dimensions is from the source value to the target value. Therefore, if we want to obtain an $(n+1)$-dimensional vector from an $n$-dimensional vector, we should add the extra components to the beginning of the vector, rather than to the end. Of course, the condition of $\nabla\cdot\mathfrak{V}_{n}$ should be satisfied anyway.

For $(\alpha, \beta, \gamma) = (3,2,1)$, we have
\begin{conjecture}
    $\lim_{n\to\infty}\mathbb{F}_n(1) =\mathbb{N} - 2\mathbb{N}$.
\end{conjecture}
\begin{proposition}
\begin{align*}
    \mathbb{V}_{n+1}(b) &= \left\{ (v, \bm{v}) \mid 2^{v}\cdot f \equiv 1\pmod{3}, v\in\mathbb{N}, f \in \mathbb{F}_{n}(b), \bm{v} \in \mathbb{V}_{n}(b) \right\},\\
    \mathbb{F}_{n+1}(b) &= \left\{ (2^{v}\cdot f  - 1)/3 \mid 2^{v}\cdot f \equiv 1\pmod{3}, v\in\mathbb{N}, f \in \mathbb{F}_{n}(b) \right\}.
\end{align*}
\end{proposition}
\begin{proposition}
    \begin{align*}
        \mathbb{V}_{n+1}(b) &= \left\{ (v, \bm{v}) \mid 2^{v}\cdot f \equiv 1\pmod{3}, v\in\mathbb{N}, f \in \mathbb{F}_{n}(b), \bm{v} \in \mathbb{V}_{n}(b) \right\} \\
        &= \bigcup_{f\in \mathbb{F}_{n}(b) - \mathbb{S}} \left\{ (2k + s, \bm{v}) \mid  k\in \mathbb{N}, \bm{v} \in \mathbb{V}_{n}(b) \right\}\\
        &= \bigcup_{f\in \mathbb{F}_{n}(b) - \mathbb{S}}\left\{ \sum_{i=0}^{(n+1)-1}(2k_i\cdot3^{i-1} + s_i)\bm{e}_{i+1} \mid  \forall k_i\in \mathbb{N} \right\},\\
        \mathbb{F}_{n+1}(b) &= \bigcup_{f\in \mathbb{F}_{n}(b) - \mathbb{S}} \left\{f_0 = (2^{s}\cdot f  - 1)/3, f_{k+1} = 4f_{k} + 1 \mid k\in\mathbb{N}\right\} \\
        &= \bigcup_{f\in \mathbb{F}_{n}(b) - \mathbb{S}} \left\{\left.\frac{1}{3}\cdot 2^{2k+s}\cdot f - \frac{1}{3}\right\vert k\in \mathbb{N}\right\}.
    \end{align*}
    Where the $\mathbb{S}$ is all primitive source defined as \cref{collatz-primitive-source}, and
    \begin{align*}
        s_i &= \min_{v\in\mathbb{N}_{\ge1}}\{2^{v}\cdot f_i \equiv 1\pmod{3} \mid f_i\in\mathbb{F}_{i}(b) - \mathbb{S} \},\\
        s &= \min_{v\in\mathbb{N}_{\ge1}}\{2^{v}\cdot f \equiv 1\pmod{3} \mid f\in\mathbb{F}_{n}(b) - \mathbb{S} \}.
    \end{align*}
\end{proposition}
This proposition provides a very important insight: the source values of a number are not entirely chaotic; they can be precisely categorized into different levels.
However, although we have provided expressions for $\mathbb{V}_{n+1}(b)$ and $\mathbb{F}_{n+1}(b)$, they are still complex and difficult to use. They are merely another equivalent expression of the original problem and are hard to use for finding counter examples. That's because they are still based on $\mathbb{F}_{n}(b)$, and therefore still require recursively representing all the first $n$ $\mathbb{F}$s to represent the $(n+1)$-th.

\section{Distribution of Multi-step Collatz Primitive Source}
Now let's consider the Collatz primitive source. They are those number that can't trace back again.
Since if $b \cdot \beta^m = \alpha\cdot a + \gamma$, then $m \in \Log_{\beta,\alpha}(b^{-1}\cdot\gamma)$.
But not all $b$ is reversible modulo $\alpha$, and not all $b^{-1} \cdot \gamma$ is a remainder of a power of $\beta$.

\begin{definition}[Collatz primitive source]~\label{collatz-primitive-source}
    For $b\in \mathbb{N}-\beta\mathbb{N}$, if there is a sequence $\bm{v}\in \mathbb{N}^{n}$ such that $f_{n}(\bm{v})$ is not reversible modulo $\alpha$, or
    \[\Log_{\beta, \alpha}(f_{n}(\bm{v})^{-1} \cdot \gamma) = \varnothing, \]
    then we say $\bm{v}$ is a $n$-step Collatz primitive source vector of $b$ and $f_{n}(a)$ is a $n$-step Collatz primitive source of $b$.
    Notice that a number $a$ is not reversible modulo $\alpha$ equivalents $\gcd(a, \alpha) > 1$. For these, we say they are trivial primitive source.
    Others we call them non-trivial primitive source.\newline
    And all non-trivial primitive source are denoted as 
    \[\tilde{\mathbb{S}} = \{b \in \mathbb{N}-\beta\mathbb{N} \mid \gcd(b, \alpha) = 1 \text{ and } \Log(b^{-1}\cdot\gamma) = \varnothing\}, \]
    all primitive source denotes as
    \[ \mathbb{S} = \{b \in \mathbb{N}-\beta\mathbb{N} \mid \gcd(b, \alpha) > 1\}\cup \tilde{\mathbb{S}}. \]
\end{definition}
If $f_n(\bm{v}) \in \mathbb{S}$, then we have $\forall k \in \mathbb{N}$, there is
\[ f_{n}(\bm{v})^{-1} \cdot \gamma \not\equiv \beta^{k} \pmod{\alpha}, \]
which equivalents
\[ \alpha^{n + 1} \nmid (\beta^{\mathcal{V}_n}\cdot b - \gamma \cdot \mathfrak{V}_n)^{-1} - \alpha^{n}\cdot\beta^{k}, \]
where the inverse is under the module of $\alpha^{n+1}$.
We can find that the distribution of Collatz primitive source deeply relay on $\beta^{k}$ module $\alpha$.
\begin{proposition}
    If $b \in \mathbb{S}$, then $\forall m \in \mathbb{Z}$, if $\beta \nmid (b + m\cdot\alpha)$, there is $b + m\cdot\alpha \in \mathbb{S}$.
\end{proposition}
But if $\beta \mid (b + m\cdot\alpha)$, then we have to turn it to $b' \equiv \frac{b + m\cdot\alpha}{\beta^{\Delta}}\pmod{\alpha}$.
This equivalence provides a new equation $b'\cdot \beta^{\Delta} = b + m'\cdot\alpha$, which is also a Collatz reduce formula, which $\gamma = b$.

\begin{proposition}
    If $\alpha \in \mathbb{P}$, then $\tilde{\mathbb{S}} = \varnothing$.
\end{proposition}

And for $(\alpha, \beta, \gamma) = (3,2,1)$, from \cref{prop:single-vary-rule} we have
\begin{proposition}
    Suppose $f_{n}(\bm{v}) \in \mathbb{S}$, then $\left\{f_{n}(\bm{v} + 6k\bm{e}_1) \mid k\in\mathbb{N}\right\} \subset \mathbb{S}$.
\end{proposition}
% !TEX root = main.tex
\section{Distribution of Collatz Primitive Source}~\label{sec:distribution-of-collatz-primitive-source}

Now let's consider the Collatz primitive source mentioned in \cref{thm:layers-of-reduce-space-and-source-value-set}. They are those number that can't trace back again.
And compaired with the total primitive source, we are more concerned about its distribution in source value set. 

\begin{definition}[Collatz primitive source]~\label{collatz-primitive-source}
    % For $b\in \widehat{\mathbb{N}}_{\beta}$, if there is a sequence $\bm{v}\in \mathbb{N}^{n}$ such that $f_{n}(\bm{v})$ is not reversible modulo $\alpha$, or
    % \[\Log_{\beta, \alpha}(f_{n}(\bm{v})^{-1} \cdot \gamma) = \varnothing, \]
    % then we say $\bm{v}$ is a $n$-step Collatz primitive source vector of $b$ and $f_{n}(\bm{v})$ is a $n$-step Collatz primitive source of $b$.
    For $b\in \widehat{\mathbb{N}}_{\beta}$, if not exists $m \ge 1$ and $a \in \widehat{\mathbb{N}}_{\beta}$ such that
    \[ b \cdot \beta^{m} = \alpha \cdot a + \gamma, \]
    then we say $b$ is a Collatz primitive source.
    All primitive source are denoted as $\mathbb{S}$.
    If $a \not\in\mathbb{S}$, we say $a$ is Collatz reversible.
\end{definition}
\begin{proposition}
    If $\alpha \mid \gamma$, then $\mathbb{S} = \{b \in \widehat{\mathbb{N}}_\beta: a\nmid b\}$.
\end{proposition}
\begin{proof}
    If $\alpha\mid \gamma$, $\forall a\in\widehat{\mathbb{N}}_\beta$, we have
    \[ a_1 \cdot \beta^{m} = \alpha \cdot a + \gamma = \alpha \cdot \left(a + \frac{\gamma}{\alpha}\right). \]
    That means $\alpha\mid a_1$. Thus if $\alpha \nmid b$, then $b$ is a primitive source.
\end{proof}
From \cref{cor:modulo-alpha-of-layers}, we have
\begin{proposition}
    If $\alpha \mid \gamma$, $\mathbb{F}_{1}(b) \not\subset \mathbb{S}$.
\end{proposition}
Suppose $\alpha \nmid \gamma$.
Since if $b \cdot \beta^m = \alpha\cdot a + \gamma$, then $m \in \Log_{\beta,\alpha}(b^{-1}\cdot\gamma)$.
But not all $b$ is reversible modulo $\alpha$, and not all $b^{-1} \cdot \gamma$ is a remainder of a power of $\beta$.
\begin{definition}
    Suppose $\alpha \nmid \gamma$. Given $b \in \mathbb{S}$. Or it is not reversible modulo $\alpha$, or
    \[\Log_{\beta, \alpha}(b^{-1} \cdot \gamma) = \varnothing. \]
    Notice that $b$ is not reversible modulo $\alpha$ equivalents $\gcd(b, \alpha) > 1$. For these, we say they are trivial primitive source.
    Others we call them non-trivial primitive source.
    And all trivial primitive source are denoted as 
    \[\hat{\mathbb{S}} = \{b \in \widehat{\mathbb{N}}_{\beta} \mid \gcd(b, \alpha) > 1\} = \bigcup_{p\in\Pi(\mathbb{P}\mid\alpha)} p\widehat{\mathbb{N}}_{{\beta}}, \]
    all non-trivial primitive source are denoted as
    \[ \tilde{\mathbb{S}} = \{b \in \widehat{\mathbb{N}}_{\beta} \mid \gcd(b, \alpha) = 1 \text{ and } \Log(b^{-1}\cdot\gamma) = \varnothing\}. \]
\end{definition}
Immidiately, we have
\begin{proposition}
    $\mathbb{S} = \tilde{\mathbb{S}} \cup \hat{\mathbb{S}}$.
\end{proposition}
\begin{proposition}~\label{prop:primitive-source-of-prime-alpha}
    If $\alpha \in \mathbb{P}$ and $\ord_{\alpha}\beta = \alpha - 1$, then $\tilde{\mathbb{S}} = \varnothing$.
\end{proposition}
Suppose $\alpha \nmid \gamma$. From \cref{cor:modulo-alpha-of-layers}, we have
\begin{proposition}
    If $\mathbb{F}_{1}(b) \cap \mathbb{S} \ne \varnothing$ and $\alpha^{2}\mid (\beta^{\Phi_1} - 1)$, then $\mathbb{F}_{1}(b) \subset \mathbb{S}$.
\end{proposition}
\begin{proposition}
    If $\mathbb{F}_{1}(b) \subset \mathbb{S}$, then $\exists d \in \Pi(\mathbb{N}\mid\alpha)$ that $d > 1$, such that $d\alpha\mid \gamma (\beta^{\Phi_1} - 1)$.
\end{proposition}
\begin{proof}
    Let $d = \gcd(P\cdot\gamma, \alpha)$.
    At first we have $A_k \equiv kP\cdot\gamma + A_0 \pmod{\alpha} \in \mathbb{S}$. If $d = 1$ then $A_k$ runs all remainder modulo $\alpha$. From \cref{thm:primitive-source-are-periodic}, we will derive that $\mathbb{N}\subset\mathbb{S}$. That's impossiable. Thus $d > 1$. Let
    \[ d = \gcd(P\cdot\gamma, \alpha) \in \Pi(\mathbb{N}\mid\alpha), \]
    we have $d \mid P\cdot\gamma = \frac{\beta^{\Phi_1} - 1}{\alpha}\cdot \gamma$, i.e. $d\alpha\mid \gamma (\beta^{\Phi_1} - 1)$.
\end{proof}
And from \cref{cor:modulo-alpha-of-layers}, we also have
\begin{theorem}~\label{thm:source-set-dijoint-primitive-source}
    Suppose $\alpha\in\mathbb{P}$, $\alpha^{2}\nmid (\beta^{\Phi_1} - 1)$ and $\alpha\nmid\gamma$.
    For $b\not\in\mathbb{S}$, we have $\forall n\ge 1$, $\mathbb{F}_{n}(b) \not\subset \mathbb{S}$ and $\mathbb{F}_{n}(b) \cap \mathbb{S} \ne \varnothing$ and $\forall a \rightsquigarrow b$, $\mathbb{F}_{1}(a)/\alpha\mathbb{Z} = \mathbb{Z}/\alpha\mathbb{Z}$.
\end{theorem}
\begin{lemma}~\label{lem:all-source-value-set-has-primitive-source}
    For $b \not\in\mathbb{S}$, if there is $\mathbb{F}_{1}(b) \not\subset \mathbb{S}$, then we have $\forall n\ge 1$, $\mathbb{F}_{n}(b) \not\subset \mathbb{S}$ and $\mathbb{F}_{n}(b) \cap \mathbb{S} \ne \varnothing$.
\end{lemma}
We can  verify that arguments $(\alpha, \beta, \gamma) = (3,2,1)$ follow the requirements of \cref{thm:source-set-dijoint-primitive-source}. This is precisely the type of Collatz system we will be discussing.

Finally, at end of this ection, we provide a theorem to describe the distribution of primitive sources in $\widehat{\mathbb{N}}_{\beta}$:
\begin{theorem}~\label{thm:primitive-source-are-periodic}
    Let $b \in \mathbb{S}$. $\forall m \in \mathbb{Z}$, if $b + m\alpha > 0$, then $\ram_{\beta}(b + m\alpha) \in \mathbb{S}$.
\end{theorem}
\begin{proof}
    Trivial for $\alpha \mid \gamma$. Suppose $\alpha \nmid \gamma$.

    For those $b \in \hat{\mathbb{S}}$, we have $\gcd(b, \alpha) > 1$. Let $d = \gcd(b, \alpha)$, then we have $\gcd(d, \beta) = 1$. Thus $d\mid \ram_{\beta}(b + m\alpha)$ means $\ram_{\beta}(b + m\alpha) \in \hat{\mathbb{S}}$.

    For those $b \in \tilde{\mathbb{S}}$, we have $\gcd(b + m\alpha,\ \alpha) = 1$. If $\ram_{\beta}(b + m\alpha) \not\in \mathbb{S}$, then $\exists k \ge 1$, such that $\beta^{k}\equiv \ram_{\beta}(b + m\alpha)^{-1} \cdot \gamma \pmod{a}$. Thus we have
    \vspace{-0.5em}
    \begin{align*}
        \beta^{k - \fac_{\beta}(b + m\alpha)} &= \beta^{k}\cdot \raf_{\beta}(b + m\alpha)^{-1} \\
        &\equiv \raf_{\beta}(b + m\alpha) \cdot \ram_{\beta}(b + m\alpha)^{-1} \cdot \gamma &&\pmod{a} \\
        &\equiv (b + m\alpha)^{-1} \cdot \gamma &&\pmod{\alpha} \\
        &\equiv b^{-1} \cdot \gamma &&\pmod{\alpha},
    \end{align*}
    which implies that $b \not\in \mathbb{S}$, conflicts with the hypothesis. Therefore, it indicates that there must be $\ram_{\beta}(b + m\alpha)\in \tilde{\mathbb{S}}$.
\end{proof}
\begin{remark}
    Notice that $\beta^{k}$ are periodic modulo $\alpha$, here we should pick a large enough $k$ to make sure $k - \fac_{\beta}(b + m\alpha) \ge 0$.
\end{remark}
In fact, not only the original source, but all source values that preserve the reversible structure (reduction sequence) follow a rule similar to \cref{thm:primitive-source-are-periodic}. This is why we discuss $\Phi_{k+1} \mid \alpha\Phi_{k}$, where these source values that preserve the reversible structure are distributed at intervals of some factor of $\alpha$. If we expand our perspective from $\mathbb{N}$ to $\mathbb{Q}$ but those denominator is a power of $\alpha$, then we can find that $b + m\alpha$ also preserves a reduction sequence centered at $b$.
\input{SECTION-densities-that-preserve-algebra-between-source-value-sets}

% \section{Continuous and Limits}\label{sec:continous-and-limits}
% To analyze the continuous and limits of Collatz sequence, we have to embed Collatz sequences in to $\mathbb{R}^{n}$.
%  \begin{definition}
%  Given $n$ dimensions vector $\bm{v} = (v_{1}, v_{2}, \cdots, v_{n}) \in \mathbb{R}^{n}$, obviously we can define its Collatz accumulation and Collatz complement by components as reduce terms.
%  With these, moreover, we can define
%      \[f(\bm{v}) = \frac{\gamma \cdot \mathfrak{A}_n}{\beta^{\mathcal{A}_n} - \alpha^n} \in \mathbb{R}  \]
%      on $\mathbb{R} - \{\mathcal{A}_n \ne n\log_{\beta}{\alpha}\}$ called Collatz value.
%  \end{definition}

%  From the \cref{theorem:any-cycle}, we have
%  \begin{conjecture}
%      If sub-sequence $a_{k_i}$ has a limit in $\mathbb{Q}_\beta$, then
%      \[ \lim_{i\to\infty}^{\mathbb{Q}_\beta} a_{k_i} = \lim_{i\to\infty}^{\mathbb{R}}\frac{\gamma \cdot \mathfrak{A}_{k_i}}{\beta^{\mathcal{A}_{k_i}} - \alpha^{k_i}}. \]
%  \end{conjecture}

\appendix
% !TEX root = main.tex
\section{Number Theory Tools}\label{sec:number-theory-tools}
To analysis the relationship between the reduces and the items, we introduced some notations from number theory.

Since we used a lot of $a\mid b$ and $a\nmid b$ but they are not in $\mathbb{Z}$, at first, we have to clarify the definition of divisibility.
\begin{definition}\label{def:extend-divisibility}
Given two numbers $\frac{p_1}{q_1}, \frac{p_2}{q_2} \in \mathbb{Q}[X]$ and they are irreducible, define $\frac{p_1}{q_1}\mid \frac{p_2}{q_2}$ as $\exists t \in \mathbb{Z}_{\ne0}$ such that $t\cdot \frac{p_1}{q_1} = p_2$. If not, we say $\frac{p_1}{q_1}\nmid \frac{p_2}{q_2}$. Of course we stipulate that $\forall \frac{p_1}{q_1} \in \mathbb{Q}[X]$, $\frac{p_1}{q_1}\mid 0$. For $\mathbb{N}$, we treat as that $q = 1$.
\end{definition}
\begin{proposition}
    $\frac{p_1}{q_1}\mid \frac{p_2}{q_2} \Longleftrightarrow p_1\mid p_2$.
\end{proposition}
\begin{proposition}
    If $a \mid c$ and $b \mid c$, and $\gcd(a, b) = 1$, then $ ab \mid c$.
\end{proposition}
Below are number theory tools we introduced in this paper and their propositions.
\begin{definition}\label{def:order-under-modulo}
    Given two number $a, b\in \mathbb{N}_{_{>1}}$ that $\gcd(a, b) = 1$, denote
    \[\ord_{a}b = \min_{x \in \mathbb{N}_{_{\ge 1}}}\left\{b^{x} \equiv 1 \pmod{a} \right\}\]
    called the discrete order of $b$ modulo $a$.
\end{definition}
\begin{proposition}
    $\ord_{a}1 = 1$.
\end{proposition}
\begin{proposition}
    If $c \mid a$, then $\ord_{c}b \mid \ord_{a}b$.
\end{proposition}
\begin{proposition}\label{prop:power-order-mod}
    $\frac{b^{s\ord_{a} b} - 1}{b^{\ord_{a} b} - 1} \equiv s \pmod{a}$.
\end{proposition}
\begin{definition}\label{def:log-under-modulo}
Given $x\in \mathbb{N}_{>1}$.
If there exists $y \in \mathbb{N}_{_{\ge 1}}$ such that $a^y \equiv x \pmod{b}$ and $y \le \ord_{b}a$, then we denote $\log_{a, b}x = y$.
If there not exists a such $y$, we will always regard $\log_{a, b}x = \infty$.
\end{definition}
\begin{proposition}
    $a^{\log_{a, b}x} \equiv x \pmod{b}$.
\end{proposition}
\begin{proposition}
    $\forall n\in\mathbb{N}$, we have $\log_{a, b}(x + nb) = \log_{a,b}x$.
\end{proposition}
\begin{proposition}
    $\forall n \in \mathbb{N}$ we have $\log_{a+nb, b}x = \log_{a, b}x$.
\end{proposition}
\begin{proposition}
    Suppose $\log_{a, b}x < \infty$ and $\log_{a, b}y < \infty$, then we have
    \[ \log_{a, b}(x \cdot y) = (\log_{a, b}x + \log_{a, b}y) \pmod{\ord_{b}a}. \]
\end{proposition}
\begin{proposition}
    Suppose $\log_{a, b}x < \infty$.
    Given $n \in \mathbb{N}$, we have
    \[ \log_{a, b}(x ^{n}) = n\log_{a, b}x \pmod{\ord_{b}a}. \]
\end{proposition}
\begin{proposition}
    Suppose $\log_{a,b}x < \infty$, then $\log_{a, b}(x^{-1}) = \ord_{b}a - \log_{a,b}x$.
\end{proposition}
\begin{definition}\label{def:log-set-under-modulo}
If $\log_{a, b}x < \infty$, then we denote that
\begin{align*}
    \Log_{a, b}x &= \ord_b(a)\mathbb{Z} + \log_{a, b}x, \\
    - \Log_{a, b}x &= \ord_b(a)\mathbb{Z} - \log_{a, b}x.
\end{align*}
If $\log_{a, b}x = \infty$, we denote $\Log_{a, b}x = \varnothing$.
\end{definition}
We call $\log_{a, b}(x)$ the principal logarithm of $x$ on $b$ modulo $a$, and $\Log_{a, b}(x)$ the logarithms of $x$ based on $b$ modulo $a$, and we will define $\log_{a, b}1 = 0$ and $\Log_{a, b}1 = \ord_{b}(a)\mathbb{Z}$.
Without special state, we will always suppose the $x > 1$ in $\log_{a, b}x$ and $\Log_{a,b}x$.

\begin{proposition}
    $\forall n\in \Log_{a,b}x$, we have $a^{n} \equiv x \pmod{b}$.
\end{proposition}
\begin{proposition}
    $\Log_{a,b}(x^{-1}) = -\Log_{a,b}x$.
\end{proposition}
\begin{definition}\label{def:generated-subgroup}
    Define $(a;b) = \{ a^{x} \mid 0 \le x < \ord_{b}(a) \}$ as the generated subgroup of $a$ modulo $b$.
\end{definition}
\begin{definition}
    Given $A\subset \mathbb{Z}$ and $a\in \mathbb{N}$. We define $A/a = \frac{A}{a} = \{ x\fmod a\mid x\in A \}$ as the remainders of $A$ modulo $a$.
\end{definition}
\begin{proposition}
    $(a^{-1};b)/b = (a;b)/b$, and $(a;b)/b \le (\mathbb{Z}/b)^{*}$ by multiplication.
\end{proposition}
% \begin{proof}
%     $a^{x} \cdot a^{y} = a^{x + y} \equiv a^{(x + y) \fmod\ord_{b}a} \pmod{b}$ and $a^{(x + y) \fmod\ord_{b}a}\in (a;b)$. Besides, $a^{-x} \equiv a^{\ord_{b} a - x} \pmod{b}$ and $a^{\ord_{b} a - x}\in (a;b)/b$. 
% \end{proof}
\begin{definition}\label{def:power-and-power-count}Define $\val_{a}b = \max_{n \in \mathbb{N}}\{n: a^n \mid b\}$ as the exponent of maximum power of $a$ in $b$ on $a$.
    Define $\pow(a\mid b) = a^{\val_{a}b}$ as the maximum power of $a$ in $b$. 
    And we define $\rel(a\mid b) = \frac{b}{\pow(a\mid b)}$ called the remainder part of $b$ with respect to power of $a$.
\end{definition}
\begin{proposition}
    If $c \mid a$ and $a \mid b$, then $\val_{c}b \ge \val_{c}a \cdot \val_{a}b$.
\end{proposition}
\begin{definition}\label{def:radical}
    Define $\rad a = \prod_{p \in \mathbb{P}}^{p \mid a}p$ as the radical of $a$. Given $S\subset\mathbb{N}$, define $\Omega(S\mid a) = \{p \in S: p \mid a\}$ as all factors of $a$ in $S$. If $S = \mathcal{P}$, we call $\Omega(\mathbb{P}\mid a)$ radical set of $a$.
\end{definition}
\begin{definition}\label{def:filtered-radical-powers}
    Define $\raf_{a}b = \prod_{p \in \Omega(\mathbb{P}\mid a)}\pow(p\mid b)$ as the filtered radical power of $b$ with respect to $a$. And we also call it the $a$-filtered radical powers of $b$ or the radical $a$-part of $b$. Specially, we define $\raf_{1}b = 1$.
    And same like to the $\rel$,  we define $\ram_{a}b = \frac{b}{\raf_{a}b}$ called the remainder part of $b$ with respect to radical $a$-part.
\end{definition}
\begin{proposition}
    $\rad a = \prod_{p \in \Omega(\mathbb{P}\mid a)}p$
\end{proposition}
\begin{proposition}\label{prop:self-containment}
    $\raf_{a}b \mid b$.
\end{proposition}
\begin{proposition}
    $\raf_{a}a = a$.
\end{proposition}
\begin{proposition}\label{prop:raf-idempotency}
    $\raf_{a}\raf_{a}b = \raf_{a}b$.
\end{proposition}
\begin{proposition}\label{prop:raf-completely-multiplicative}
    $\raf_{a}(b_1b_2) = \raf_{a}b_1\cdot\raf_{a}b_2$.
\end{proposition}
\begin{proposition}\label{prop:raf-inclusion-exclusion-principle}
    $\raf_{a_1 a_2}b = \frac{\raf_{a_1}b \cdot \raf_{a_2}b}{\raf_{\gcd(a_1, a_2)}b}$.
\end{proposition}
\begin{proposition}\label{prop:raf-power-homogeneity}
    $\raf_{a}(b^n) = (\raf_{a}b)^{n}$.
\end{proposition}
\begin{proposition}\label{prop:raf-index-monotonicity}
    $a_1 \mid a_2 \implies \raf_{a_1}b \mid \raf_{a_2}b$.
\end{proposition}
\begin{proposition}\label{prop:raf-valuation-monotonicity}
    $\val_{p}b_1 \le \val_{p}b_2 \Longleftrightarrow \raf_{p}b_1 \mid \raf_{p}b_2$.
\end{proposition}
\begin{proposition}\label{prop:raf-radical-invariance}
    $\rad a_1 = \rad a_2 \implies \raf_{a_1}b = \raf_{a_2}b$.
\end{proposition}
\begin{proposition}\label{prop:raf-gcd-distributivity}
    $\raf_{a}\gcd(b_1, b_2) = \min\{\raf_{a}b_1, \raf_{a}b_2\}$.
\end{proposition}
\begin{proposition}\label{prop:raf-lcm-distributivity}
    $\raf_{a}\lcm(b_1, b_2) = \max\{\raf_{a}b_1, \raf_{a}b_2\}$.
\end{proposition}
\begin{proposition}\label{prop:raf-limit-power-gcd-raf-equvalent}
    $\lim_{n\to\infty}\gcd(a^{n}, b) = \raf_{a}b$.
\end{proposition}
\begin{proposition}
    $\log \raf_{a}b = \sum_{p\in\mathbb{P}}^{p \mid a} \val_{p}b \log p$
\end{proposition}
\begin{proposition}\label{prop:rem-coprimality}
    $\gcd(\ram_{a}b, a) = 1$.
\end{proposition}
\begin{proposition}\label{prop:raf-ord-non-self-divisibility}
    If $p \in \mathbb{P}$, $p\nmid b$, then $\raf_{p}\ord_{p}b = 1$.
\end{proposition}
\begin{theorem}\label{thm:lte-radical-theorem}
    If $p \in \mathbb{P}$ and $p\mid (b - 1)$, then
    \[\raf_{p}(b^{n} - 1) = \left\{ \begin{aligned}
        &\raf_{p}(b - 1) \cdot \raf_{p}n, & p > 2, \\
        &\raf_{2}(b - 1) \cdot \raf_{2}n + \raf_{2}(b + 1) - 1, & p = 2.
    \end{aligned} \right.\]
\end{theorem}

\begin{definition}\label{def:radicable-order}
    $\rord_{a} b = \lcm_{p\in\mathbb{P}}^{p\mid a} \{\ord_{p} b\}$ be the radicable order.
\end{definition}

\begin{definition}\label{def:radical-height-and-width}
    Given $a \in \mathbb{N}_{\ge 1}$. Denote $\nu(a) = \max\{ \val_{p} a \mid p\in\mathbb{P}\}$ and $\omega(a) = \card\Omega(\mathbb{P}\mid a)$ called the radical height and the radical width of $a$.
\end{definition}

% Now let's introduce a kind of count function on infinite subsets of $\mathbb{N}$ from analytic number theory. Given $A \subset N$ that $\card{A} = \infty$ and $\forall a\in A$, $a > 1$.
% \begin{definition}
%     Define $\pi(A;x) = \card\{ A \le x \}$ that $x\in\mathbb{R}_{>1}$. And we define the Perron form
%     \[\widetilde{\pi}(A;x) = \left\{\begin{aligned}
%     &\pi(A;x), && x\not\in A, \\
%     &\frac{1}{2}\left[\pi(A;x^{-}) + \pi(A;x^{+})\right], && x\in A.
% \end{aligned}\right.\]
% \end{definition}
% \begin{proposition}
%     $\sum_{t\in A}^{t\le x}f(t) = \int_{1}^{x}f(t)\rmd{\pi(A;t)}$.
% \end{proposition}
% \begin{proposition}
%     $\pi(A\cup B;x) = \pi(A;x) + \pi(B; x) - \pi(A\cap B; x)$. So does for $\widetilde{\pi}(A;x)$.
% \end{proposition}
% \begin{definition}
%     Define $D(A;s) = \sum_{x\in A}x^{-s}$ as the Dirichlet series of $A$.
% \end{definition}
% \begin{proposition}
%     $D(aA;s) = a^{-s}D(A;s)$.
% \end{proposition}
% \begin{proposition}
%     If $A\otimes B$ is uniqe decomposed, then we have
%     \[ D(A \otimes B; s) = D(A;s)\cdot D(B;s). \]
% \end{proposition}
% \begin{theorem}\label{thm:count-function-from-dirichlet-series}
%     For large enough $c > 0$, we have
%     \[ \widetilde{\pi}(A;x) = \frac{1}{2\pi\rmi}\int_{c-\rmi\infty}^{c+\rmi\infty}D(A;s)\frac{x^{s}}{s}\rmd{s}. \]
% \end{theorem}
% \begin{lemma}[Perron's formula]
% For $y > 0$, we have
% \[ P(x) = \frac{1}{2\pi\rmi}\int_{c-\rmi\infty}^{c+\rmi\infty}\frac{x^{s}}{s}\rmd s = \left\{\begin{aligned}
%     &1, &&y > 1,\\
%     &1/2, &&y = 1,\\
%     &0, &&y < 1.
% \end{aligned}\right. \]
% \end{lemma}
% \begin{proof}[Proof of \cref{thm:count-function-from-dirichlet-series}]
%     \begin{align*}
%         &\frac{1}{2\pi\rmi}\int_{c-\rmi\infty}^{c+\rmi\infty}D(A;s)\frac{x^{s}}{s}\rmd s \\
%         =&\frac{1}{2\pi\rmi}\int_{c-\rmi\infty}^{c+\rmi\infty}\sum_{t\in A}t^{-s}\frac{x^{s}}{s}\rmd s \\
%         =&\sum_{t\in A}\frac{1}{2\pi\rmi}\int_{c-\rmi\infty}^{c+\rmi\infty}\left(\frac{x}{t}\right)^{s}\frac{1}{s}\rmd s \\
%         =&\sum_{t\in A}P(x/t) \\
%         =&\widetilde{\pi}(A;x).
%     \end{align*}
% \end{proof}
% \begin{definition}
%     Define $\Pi(A;x) = \sum_{m=1}^{\infty}\frac{1}{m}\pi(A;x^{1/m})$ and also the Perron form $\widetilde{\Pi}(A;x) = \sum_{m=1}^{\infty}\frac{1}{m}\widetilde{\pi}(A;x^{1/m})$.
% \end{definition}
% \begin{definition}
%     For $s > 0$, we define $L(A;s) = \prod_{x\in A}(1-x^{-s})^{-1}$.
% \end{definition}
% \begin{proposition}
%     $\ln L(A;s) = \sum_{m=1}^{\infty}\frac{1}{m}\sum_{x\in A}x^{-sm} = s\int_{1}^{\infty}\frac{\Pi(A;x)}{x^{s+1}}\rmd x$.
% \end{proposition}
% \begin{proof}
% At first, we have
% \[ \ln L(A;s) = -\sum_{x\in A}\ln(1 - x^{-s}) = \sum_{x\in A}-\ln(1 - x^{-s}). \]
% From Maclaurin's formula, $- \ln (1 - z) = \sum_{m=1}^{\infty}\frac{x^{m}}{m}$, we have
% \begin{align*}
%     \ln L(A;s) &= \sum_{x\in A}\sum_{m=1}^{\infty}\frac{x^{-sm}}{m} = \sum_{m=1}^{\infty}\frac{1}{m}\sum_{x\in A}x^{-sm} \\
%     &= \sum_{m=1}^{\infty}\frac{1}{m}\int_{1}^{\infty}x^{-sm}\rmd\pi(A;x) \\
%     &= \sum_{m=1}^{\infty}\frac{1}{m}\int_{1}^{\infty}x^{-s}\rmd\pi(A;x^{1/m}) \\
%     &= \int_{1}^{\infty}x^{-s}\rmd\left(\sum_{m=1}^{\infty}\frac{1}{m}\pi(A;x^{1/m})\right) \\
%     &= \int_{1}^{\infty}x^{-s}\rmd\Pi(A;x)
% \end{align*}
% According the rule of integration by parts, we have
% \begin{align*}
%     \ln L(A;s) &= \left[x^{-s}\Pi(A;x)\right]\bigg\vert_{1}^{\infty} - \int_{1}^{\infty}\Pi(A;x)\rmd x^{-s} = s\int_{1}^{\infty}\frac{\Pi(A;x)}{x^{s+1}}\rmd x.
% \end{align*}
% \end{proof}
% \begin{theorem}
%     Let $\mu(m)$ be the Möbius function, for $c > 0$ large enough, we have
%     \[ \widetilde{\pi}(A;x) = \frac{1}{2\pi\rmi}\sum_{m=1}^{\infty}\frac{\mu(m)}{m}\int_{c-\rmi\infty}^{c+\rmi\infty}\ln{L(A;ms)}\frac{x^{s}}{s}\rmd s. \]
% \end{theorem}
% \begin{proof}
%     Since we have Perron's formula, we are supported to have
% \begin{align*}
%     &\frac{1}{2\pi\rmi}\int_{c-\rmi\infty}^{c+\rmi\infty}\ln{L(A;s)}\frac{x^{s}}{s}\rmd s \\
%     =&\frac{1}{2\pi\rmi}\int_{c-\rmi\infty}^{c+\rmi\infty}\sum_{t\in A}\sum_{m=1}^{\infty}\frac{t^{-sm}}{m}\left(\frac{x^{s}}{s}\right)\rmd s \\
%     =&\sum_{m=1}^{\infty}\sum_{t\in A}\frac{1}{m}\left[ \frac{1}{2\pi\rmi}\int_{c-\rmi\infty}^{c+\rmi\infty}\left(\frac{x}{t^{m}}\right)^{s}\frac{1}{s}\rmd s \right] \\
%     =&\sum_{m=1}^{\infty}\sum_{t\in A}\frac{1}{m}P(x/t^{m}) \\
%     % =&\sum_{m=1}^{\infty}\frac{1}{m}\card\{t\in A \mid \frac{x}{t^{m}} > 1 \} \\
%     % =&\sum_{m=1}^{\infty}\frac{1}{m}\card\{t\in A \mid t^{m} < x \} \\
%     =&\sum_{m=1}^{\infty}\frac{1}{m}\widetilde{\pi}(A;x^{1/m}) \\
%     =&\widetilde{\Pi}(A;x).
% \end{align*}
% Using the Möbius inversion, we have
% \begin{align*}
%     \widetilde{\pi}(A;x) &= \sum_{m=1}^{\infty}\frac{\mu(m)}{m}\widetilde{\Pi}(A;x^{1/m}) \\
%     &= \sum_{m=1}^{\infty}\frac{\mu(m)}{m}\left[ \frac{1}{2\pi\rmi}\int_{c-\rmi\infty}^{c+\rmi\infty}\ln{L(A;s)}\frac{x^{s/m}}{s}\rmd s \right] \\
%     &= \frac{1}{2\pi\rmi}\sum_{m=1}^{\infty}\frac{\mu(m)}{m}\int_{c-\rmi\infty}^{c+\rmi\infty}\ln{L(A;ms)}\frac{x^{s}}{s}\rmd s.
% \end{align*}
% \end{proof}
% !TEX root = main.tex
\section{A Kind of Derivation on Groupoid}\label{sec:a-kind-of-derivation-on-groupoid}
In \cref{definition:collatz-complement} we defined the Collatz complement, they follow the proposition \cref{proposition:collatz-complement-basic-prop2}.
Now let's forget those definitions before.
We will introduce a new kind of form to subsume those properties above. It's the extend of $(\sigma, \tau)$-derivation on groupoids.

Let $\mathbb{G}(+, G)$ be a groupoid and $\mathbb{R}(+, \cdot, R)$ be a ring.
We can define two kind of functions that from $\mathbb{G}$ to $\mathbb{R}$.
\begin{definition}[groupoid homomorphism]
    Function $\mathcal{A}: \mathbb{G} \to \mathbb{R}$ is called a groupoid homomorphism, if $\forall a, b \in \mathbb{G}$ that $ab \in \mathbb{G}$, there is $\mathcal{A}(a+b) = \mathcal{A}(a) \cdot \mathcal{A}(b)$.
\end{definition}
\begin{proposition}
    If $\exists a \in \mathbb{G}$ that $\mathcal{A}(a) = 0$, then $\forall b \in \mathbb{G}$, we have $\mathcal{A}(b) = 0$.
    We called that $\mathcal{A}$ is a zero groupoid homomorphism, and denote $\mathcal{A} = 0$.
    Otherwise, we say $\mathcal{A} \ne 0$.
\end{proposition}
\begin{proposition}
    If $\mathcal{A} \ne 0$, then $\forall a \in \mathbb{G}$, there is $\mathcal{A}(a)\cdot\mathcal{A}(-a) = \mathcal{A}(0) = 1$. In other words, $\mathcal{A}(-a) = \mathcal{A}(a)^{-1}$. This requires the ring $\mathbb{R}$ at least is a division ring.
\end{proposition}
\begin{proposition}
    Denote $\mathrm{Hom}(\mathbb{G},\mathbb{R})$ as the set of all groupoid homomorphisms on $\mathbb{R}$, then $\mathrm{Hom}(\mathbb{G},\mathbb{R})$ is a group by
    \[ (\mathcal{A}\cdot\mathcal{B})(a) = \mathcal{A}(a)\cdot\mathcal{B}(a). \]
\end{proposition}
\begin{definition}[($\mathcal{A}$, $\mathcal{B}$)-derivation]~\label{def:quasi-derivation-on-groupoid}
    Given two groupoid homomorphisms $\mathcal{A}$ and $\mathcal{B}$.
    There's another function $\mathfrak{L}: \mathbb{G} \to \mathbb{R}$ such that $\forall a, b \in \mathbb{G}$, if $a + b \in \mathbb{G}$, then
    \[ \mathfrak{L}(a+b) = \mathcal{A}(a)\cdot\mathfrak{R}(b) + \mathfrak{R}(a)\cdot\mathcal{B}(b) \]
    called the left ($\mathcal{A}$, $\mathcal{B}$)-derivation.
    And function $\mathfrak{R}: \mathbb{G} \to \mathbb{R}$
    \[ \mathfrak{R}(a+b) = \mathcal{A}(b)\cdot\mathfrak{R}(a) + \mathfrak{R}(b)\cdot\mathcal{B}(a) \]
    called the right ($\mathcal{A}$, $\mathcal{B}$)-derivation.
\end{definition}
Actually, ($\mathcal{A}$, $\mathcal{B}$)-derivation is the extension of the $(\sigma, \tau)$-derivation on groupoid.
Given a ($\mathcal{A}$, $\mathcal{B}$)-derivation $\mathfrak{U}$, left or right both have
\begin{proposition}
    $\mathfrak{U}(0) = 0$.
\end{proposition}
\begin{proposition}
    $\mathcal{A}(a)\cdot\mathfrak{U}(-a) + \mathfrak{U}(a)\cdot\mathcal{B}(-a) = 0$.
\end{proposition}

Denote $\partial_{R}(\mathcal{A}, \mathcal{B})$ as the set of all right ($\mathcal{A}$, $\mathcal{B}$)-derivations and $\partial_{L}(\mathcal{A}, \mathcal{B})$ as the set of all left ($\mathcal{A}$, $\mathcal{B}$)-derivations.
If the left or right is not matter, i.e. $\mathbb{R}$ is commutative, we just denote them $\partial_{*}(\mathcal{A}, \mathcal{B})$.
Our discussion of ($\mathcal{A}$, $\mathcal{B}$)-derivations will base on $\mathbb{R}$ is commutative.
\begin{proposition}
    $\partial_{*}(\mathcal{A}, \mathcal{B})$ is a $\mathbb{R}$-module by
    \[ (s\mathfrak{U} + t\mathfrak{V})(a) = s\mathfrak{U}(a) + t\mathfrak{V}(a). \]
\end{proposition}
Also, we denote $\mathfrak{U} = 0$ if it is a zero ($\mathcal{A}$, $\mathcal{B}$)-derivation.
However, we need another restriction of ($\mathcal{A}$, $\mathcal{B}$)-derivations to describe our problem.

We denote $\mathbb{P}(\mathbb{G}) = \{ p\in\mathbb{G} \mid p \text{ is inseperable}\}$. With this we have
\begin{definition}[kernel of ($\mathcal{A}$, $\mathcal{B}$)-derivation]
    We say $\mathfrak{U}(a)$ is a kernel on $a$ of $\mathfrak{U}$ if $a \in \mathbb{P}(\mathbb{G})$. If $\mathbb{P}(\mathbb{G}) = \varnothing$, we define all kernels of $\mathfrak{U}$ are same.
\end{definition}
\begin{definition}
    If $\exists t \in\mathbb{N}_{\ge 1}$ such that $\forall a_{k}\in \mathbb{P}(\mathbb{G})$, such that $F:\mathbb{G} \to \mathbb{R}$ satisfied $F(a_{k + t}) = F(a_k)$, then we say $F$ is a cycle kernel form with $t$ is a period of it.
\end{definition}
All cycle kernel ($\mathcal{A}$, $\mathcal{B}$)-derivations with period $t$ are denoted as $\partial_{\circ(t)}(\mathcal{A}, \mathcal{B})$.
\begin{definition}
    Given a ($\mathcal{A}$, $\mathcal{B}$)-derivation $\mathfrak{U} \in \partial_{*}(\mathcal{A}, \mathcal{B})$, if $\forall p, q\in \mathbb{P}(\mathbb{G})$, there is $\mathfrak{U}(p) = \mathfrak{U}(q)$, we say $\mathfrak{U}$ is kernel independent.
\end{definition}
All kernel independent ($\mathcal{A}$, $\mathcal{B}$)-derivations are denoted as $\partial(\mathcal{A}, \mathcal{B})$.
Of course, we have
\begin{proposition}
    $\partial(\mathcal{A}, \mathcal{B}) = \partial_{\circ(1)}(\mathcal{A}, \mathcal{B})$.
\end{proposition}
\begin{theorem}
    Suppose $\mathcal{A}, \mathcal{B} \in\mathrm{Hom}(\mathbb{G}, \mathbb{R})$.
    If $c = \sum_{i = 0}^{n-1} a_i \in \mathbb{G}$, denote $P_k = \sum_{i = 0}^{k-1}a_i$ and $S_{k} = \sum_{i = k + 1}^{n-1}a_i$, then
    $\forall \mathfrak{U} \in \partial_{*}(\mathcal{A}, \mathcal{B})$ we have
    \[ \mathfrak{U}(c) = \sum_{k = 0}^{n-1}\mathcal{A}(P_k)\cdot\mathfrak{U}(a_k)\cdot\mathcal{B}(S_k). \]
\end{theorem}
\begin{corollary}
    If $\mathcal{A}$ and $\mathcal{B}$ are periodic that peroid is $t$, then $\forall \mathfrak{U} \in \partial_{\circ(t)}(\mathcal{A}, \mathcal{B})$, we have
    \[ \mathfrak{U}(\sum_{i = t}^{n+t-1} a_i) = \mathfrak{U}(\sum_{i = 0}^{n - 1} a_i)  \]
\end{corollary}
Specially, For the kernel independent ($\mathcal{A}$, $\mathcal{B}$)-derivation, we have
\begin{corollary}
    If $\mathfrak{U} \in \partial(\mathcal{A}, \mathcal{B})$ that $\mathfrak{U}(p) = \gamma$, then
    \[ \mathfrak{U}(c) = \gamma \cdot\sum_{k = 0}^{n-1}\mathcal{A}(P_k)\cdot\mathcal{B}(S_k). \]
\end{corollary}

Finially back to the Collatz complement. With arguments $(\alpha, \beta, \gamma)$, all Collatz complement $\mathfrak{A} \in \partial(\beta^{\mathcal{A}_{*}}, \alpha^{*})$, where the groupoid is $[\mathbb{Z}\times\mathbb{Z}]_{\rightsquigarrow}$ that $i \rightsquigarrow j$ as its element.
Inseparable elements of it is $\mathbb{P}([\mathbb{Z}\times\mathbb{Z}]_{\rightsquigarrow}) = \{i \rightsquigarrow (i + 1) \mid i \in \mathbb{Z}\}$. And the kernel $\mathfrak{A}(i \rightsquigarrow i + 1) = \mathfrak{A}_{i+1}^{i} = \gamma$.
Besides, the Collatz complement of Conway's extension of Collatz problem is actually those $\partial_{*}(\mathcal{A}, \mathcal{B})$, they are not kernel independent, but its kernels are $\beta$-peroidic. Hence, no matter our Collatz complement, or Collatz complement of Conway's extension, just let $\mathcal{A}$ and $\mathcal{B}$ are $\beta$-periodic, they will all be peroidic.


\bibliography{Collatz}
\bibliographystyle{unsrt}
\nocite{*}

\end{document}
